\documentclass[11pt]{amsart}

\usepackage[utf8]{inputenc}
\usepackage[T1]{fontenc}
\usepackage{amsmath,amssymb,amsthm}
\usepackage{mathtools}
\usepackage{hyperref}
\usepackage[margin=1.1in]{geometry}
\usepackage{enumitem}
\usepackage{xcolor}

\hypersetup{
  colorlinks=true,
  linkcolor=blue!60!black,
  citecolor=green!50!black,
  urlcolor=blue!60!black,
}

\newtheorem{theorem}{Theorem}[section]
\newtheorem{lemma}[theorem]{Lemma}
\newtheorem{proposition}[theorem]{Proposition}
\newtheorem{corollary}[theorem]{Corollary}
\theoremstyle{definition}
\newtheorem{definition}[theorem]{Definition}
\newtheorem{remark}[theorem]{Remark}
\newtheorem{notation}[theorem]{Notation}

\newcommand{\RR}{\mathbb{R}}
\newcommand{\NN}{\mathbb{N}}
\newcommand{\PP}{\mathbb{P}}     % Leray projection
\newcommand{\BB}{\mathcal{B}}    % Bogovskii
\newcommand{\EE}{\mathcal{E}}    % weighted energy
\newcommand{\DD}{D}              % annular gradient norm
\newcommand{\Lweak}[1]{L^{3,\infty}(\RR^{#1})}
\newcommand{\norm}[1]{\|#1\|}
\newcommand{\abs}[1]{|#1|}
\newcommand{\loc}{\mathrm{loc}}
\newcommand{\supp}{\mathrm{supp}}

\DeclareMathOperator{\dv}{div}
\DeclareMathOperator{\curl}{curl}

\title{A Liouville theorem for stationary Navier--Stokes in $L^{3,\infty}(\RR^3)$, and the endpoint gap in the Clay regularity chain}

\author{Airo Nahiru}
\address{Independent researcher}
\email{aironahiru@gmail.com}

\date{\today}

\begin{document}
\begin{abstract}
We prove that every smooth stationary divergence-free solution of
the incompressible three-dimensional Navier--Stokes equations that
lies in the Lorentz space $L^{3,\infty}(\RR^3)$ is identically
zero. The proof is entirely first-principles: a hat-function
weighted energy on fixed-ratio annuli is combined with exact
pressure elimination via the Bogovski\u{\i} operator, a Lorentz
layer-cake embedding, and a rescaled Gagliardo--Nirenberg
interpolation. Exact $\rho$-cancellation at scaling exponent
$3a/p + b = 3$ yields a superlinear differential inequality
$\EE'(\rho) \geq c\rho^{-\alpha}\EE(\rho)^\beta$ whose exponents
satisfy $\beta > 1$, $\alpha < 1$ for every $p \in (3/2, 3)$;
integration against bounded left-hand side produces a contradiction
that forces finite Dirichlet energy, which a Galdi-style testing
argument reduces to $v = 0$.

The theorem completes the Koch--Nadirashvili--Seregin--\v{S}ver\'ak
(KNSS) program for the general (non-axisymmetric) three-dimensional
case. Applied within the KNSS blow-up framework, it reduces the
full Clay Millennium regularity problem to a single structural
inheritance step: the rescaled ancient solution's uniform
$L^{3,\infty}$ bound. That step is equivalent to Phuc's 2015 open
endpoint, which remains unresolved for non-axisymmetric flows.
We give an honest account of this gap, survey the partial
closures (axisymmetric, smallness, finite singular set,
log-improvement), and describe a concrete weighted-moment
proposal. The companion corpus at
\url{https://github.com/hirunation/finish_navier_stokes} (archived at
Zenodo, DOI forthcoming) contains the full provenance-tracked
knowledge graph, nine canonical obstruction theorems, thirty-five
catalogued dead ends, and the Lean~4 formalization.
\end{abstract}

\maketitle

\tableofcontents

\section{Introduction}\label{sec:intro}

\subsection{The problem.}
For the incompressible Navier--Stokes equations in $\RR^3$,
\begin{align}\label{eq:NS}
\partial_t u - \nu\,\Delta u + (u \cdot \nabla) u + \nabla p &= 0,
& \nabla \cdot u &= 0, & u(\cdot, 0) &= u_0,
\end{align}
the Clay Millennium Prize problem asks whether every smooth
Schwartz-class divergence-free $u_0$ admits a global smooth
solution. Leray~\cite{Leray} proved existence of weak solutions
in 1934; regularity for general smooth data remains open.

The Koch--Nadirashvili--Seregin--\v{S}ver\'ak blow-up framework
\cite{KNSS} reduces the regularity question to a Liouville-type
theorem. If a smooth solution develops a finite-time singularity
at $T^*$, a rescaling procedure (\cite[\S6]{KNSS}) extracts a
bounded smooth ancient solution $v \colon \RR^3 \times
(-\infty, 0] \to \RR^3$ with $\abs{v(0, 0)} = 1$. The backward-time
limit of this ancient solution is a smooth stationary
divergence-free Navier--Stokes solution; if that limit were forced
to be identically zero, Escauriaza--Seregin--\v{S}ver\'ak backward
uniqueness~\cite{ESS03RMS} would contradict the normalization
$\abs{v(0, 0)} = 1$, closing the regularity problem.

KNSS~\cite[\S5]{KNSS} established such a Liouville theorem under
axial symmetry, using stream-function arguments that reduce 3D to
2D. For general (non-axisymmetric) three-dimensional flows their
method does not extend, and they flag this as an open problem.
The present paper supplies the general-case Liouville theorem
under the critical Lorentz-space hypothesis
$v \in L^{3,\infty}(\RR^3)$.

\subsection{Main theorem.}\label{ssec:main}

\begin{theorem}[Stationary $L^{3,\infty}$ Liouville]\label{thm:main}
Let $v \colon \RR^3 \to \RR^3$ be smooth, divergence-free, and
satisfy the stationary Navier--Stokes equations
\[
-\nu\,\Delta v + (v \cdot \nabla) v + \nabla p = 0,
\qquad \nabla \cdot v = 0
\quad \text{on } \RR^3,
\]
with $\norm{v}_{\Lweak{3}} < \infty$. Then $v \equiv 0$.
\end{theorem}

Theorem~\ref{thm:main} completes the KNSS program for the general
three-dimensional case. Its proof is self-contained and uses only
first-principles tools (the Bogovski\u{\i} operator for pressure
elimination, Lorentz layer-cake embedding, Gagliardo--Nirenberg
interpolation); no ancient-Liouville conjectures are invoked.

The critical exponents at the natural choice $p = 2$ are
$\beta = 4/3$, $\alpha = 2/3$, and $\gamma = -1/4$. The scale
exponent $\gamma$ is defined in
Theorem~\ref{thm:main-estimate} (\S\ref{sec:gn}), and the ODE
exponents $\alpha$ and $\beta$ in
Proposition~\ref{prop:ODE} (\S\ref{sec:ode}); Sections
\ref{sec:pressure}--\ref{sec:vzero} carry the full integrated
argument.

\subsection{The Clay chain and the $L^{3,\infty}$ inheritance gap.}
\label{ssec:clay-gap}
Applied within the KNSS framework, Theorem~\ref{thm:main} reduces
the Clay regularity question to a single structural step. The
rescaled ancient solution produced by \cite[\S6]{KNSS} is bounded
in $L^\infty$; what is required in addition, for
Theorem~\ref{thm:main} to apply to its backward-time limit, is a
uniform $L^{3,\infty}$ bound on the ancient solution. Scale
invariance reduces this question to the spatial decay of the
original solution in the exterior of the hypothetical singular
set, uniformly on $[0, T^*)$.

This inheritance step is currently open. It is equivalent in
scope to Phuc's 2015 open problem on the $L^{3,\infty}$ endpoint
regularity criterion~\cite[p.~4]{Phuc2015}:
\begin{quote}
``To the best of our knowledge, a criterion of local regularity
for the Navier--Stokes equations in $L^\infty(-1, 0;
L^{3,\infty}(B_1))$ is still unknown.''
\end{quote}
The obstruction is structural: the Lorentz quasi-norm
$\norm{\cdot}_{\Lweak{3}}$ is not absolutely continuous, so
exterior-tail compactness arguments that close the $L^{3,q}$
endpoint for $q < \infty$ fail at $q = \infty$. Partial closures
exist: the axisymmetric case (O\.{z}a\'nski--Palasek~\cite{OzanskiPalasek}
via Nazarov--Ural'tseva Harnack~\cite{NazarovUraltseva} on
$rv_\theta$), the smallness regime (Luo--Tsai~\cite{LuoTsai},
Kim--Kozono~\cite{KimKozono}), and partial regularity in the sense
of finite singular sets (Seregin~\cite{Seregin2019}).

Section~\ref{sec:clay-chain} describes these partial closures and
gives a concrete proposal (\S\ref{ssec:weighted-moment}) for
closing the general case under the weighted hypothesis
$(1 + \abs{x}) u_0 \in L^\infty$, which is automatic for Schwartz
data. A full proof of that proposal is outside the scope of the
present paper.

\subsection{Corollaries.}

\begin{corollary}[Four-dimensional Liouville]\label{cor:4d}
Let $v \colon \RR^4 \to \RR^4$ be smooth, divergence-free, and
satisfy the stationary Navier--Stokes equations on $\RR^4$ with
$v \in L^{4,\infty}(\RR^4)$. Then $v \equiv 0$.
\end{corollary}

Corollary~\ref{cor:4d} follows from a four-dimensional
specialization of the proof at $p = 5/2$, where the dimensional
discriminant $\gamma(n, p) = -(n-4)(n-2p)/(2n - p(n-2))$ vanishes
identically, producing the borderline logarithmic divergence
$\int \rho^{-1}\,d\rho$ against a bounded left-hand side. The
discriminant becomes positive for $n \geq 5$, and the method
provably fails there: $\beta > 1$ requires $p > n/2$ (Sobolev)
while $\alpha \leq 1$ requires $p \leq n/2$ (integration), an
incompatibility.

\begin{corollary}[Weak-solution extension]\label{cor:weak}
The conclusion of Theorem~\ref{thm:main} holds for
$v \in L^{3,\infty}(\RR^3) \cap W^{1,2}_\loc(\RR^3)$.
\end{corollary}

The $W^{1,2}_\loc$ threshold is sharp: Landau solutions
$v(x) = c/\abs{x}$ are smooth outside the origin, lie exactly in
$L^{3,\infty}(\RR^3)$, and belong to $W^{1,2}_\loc$ away from the
origin but not at the origin. They realize the boundary case.

\begin{corollary}[Quantitative bound]\label{cor:quant}
From the superlinear ODE at $p = 2$, the weighted Dirichlet
energy on a scale-$\rho_0$ annulus satisfies the explicit bound
$\EE(\rho_0) \leq C\,M^6\,\rho_0^{-1}$, with explicit constants.
\end{corollary}

\subsection{Companion corpus.}\label{ssec:companion}
The complete corpus --- $111$ provenance-tracked smart-data
entries spanning insights, lemmas, definitions, obstruction
theorems, classical axioms, anti-knowledge entries, attribution
corrections, master proof chains, and cross-cutting narrative
threads --- is maintained publicly at
\begin{center}
\url{https://github.com/hirunation/finish_navier_stokes}
\end{center}
and archived at Zenodo (DOI forthcoming). The repository includes
the Lean~4 formalization~(the \texttt{GUNS} subdirectory and Tree
A at \texttt{navier\_stokes\_official/lean\_formalization/}),
literature reviews for the nine obstruction theorems of
\S\ref{sec:obstructions}, the thirty-five catalogued dead ends,
and the desk-level attack analyses referenced throughout this
paper. Readers seeking the full architecture of the Clay-chain
program, the obstruction landscape, or the alternative-route
research programs (Hull-Ergodic, Icosahedral, Ancient Extension)
are directed there.

\subsection{Organization.}

Section~\ref{sec:notation} fixes notation. Sections
\ref{sec:hat}--\ref{sec:vzero} prove Theorem~\ref{thm:main} in
eight steps. Section~\ref{sec:extensions} proves the
Corollaries~\ref{cor:4d}--\ref{cor:quant}. Section
\ref{sec:clay-chain} describes the application to the Clay chain
and documents the $L^{3,\infty}$ inheritance gap. Section
\ref{sec:obstructions} gives a brief summary of the
nine-obstruction landscape. Section~\ref{sec:discussion}
concludes with discussion and open problems.

\section{Notation and preliminaries}\label{sec:notation}

Throughout, $\nu > 0$ is a fixed viscosity and $v$ denotes a
smooth divergence-free stationary solution of~\eqref{eq:NS} on
$\RR^3$ with
\[
M := \norm{v}_{\Lweak{3}} < \infty.
\]
We assume $M > 0$; otherwise $v = 0$ almost everywhere, hence
identically by smoothness, and there is nothing to prove.

\begin{notation}[Geometric objects]\label{not:geometry}
Fix $\theta \in (0, 1)$. For $\rho > 0$, write
\[
B_\rho := \{x \in \RR^3 : \abs{x} < \rho\},\qquad
A(\theta\rho, \rho) := B_\rho \setminus \overline{B_{\theta\rho}}.
\]
Define the \emph{hat function} $\varphi \colon [0, \infty) \to
[0, 1]$ by
\[
\varphi(t) = \begin{cases}
1 & 0 \leq t < \theta, \\
\dfrac{1 - t}{1 - \theta} & \theta \leq t < 1, \\
0 & t \geq 1.
\end{cases}
\]
The \emph{weighted energy at scale} $\rho$ is
\[
\EE(\rho) := \int_{\RR^3} \abs{\nabla v(x)}^2\, \varphi(\abs{x}/\rho)\, dx,
\]
and the \emph{annular gradient norm} is
\[
\DD(\rho) := \norm{\nabla v}_{L^2(A(\theta\rho, \rho))}^2.
\]
\end{notation}

The hat function is piecewise linear on $[\theta, 1]$ and zero
outside $[0, 1]$, smooth on the open pieces with a bounded weak
gradient. It is $C^{0,1}$ globally, and for the arguments below
it may be replaced by any fixed mollification supported in
$A(\theta\rho, \rho)$ with the same qualitative properties; we
use the piecewise-linear form throughout for the cleanest
derivative formula (cf.\ Remark~\ref{rem:hat}).

\begin{lemma}[Hat-function derivative formula]\label{lem:hat-deriv}
For every $\rho > 0$,
\begin{equation}\label{eq:hat-deriv}
\EE'(\rho)
= \frac{1}{(1-\theta)\rho^2} \int_{A(\theta\rho, \rho)} \abs{\nabla v(x)}^2 \abs{x}\, dx
\geq \frac{\theta}{(1 - \theta)\rho}\, \DD(\rho).
\end{equation}
In particular, $\EE$ is absolutely continuous and non-decreasing
in $\rho$, and
\begin{equation}\label{eq:D-bounded}
\DD(\rho) \leq \frac{1 - \theta}{\theta} \rho\, \EE'(\rho).
\end{equation}
\end{lemma}

\begin{proof}
Differentiating under the integral sign (valid by $v$ smooth and
$\varphi$ weakly differentiable with compact support in $\rho$),
\[
\EE'(\rho) = \int \abs{\nabla v(x)}^2 \frac{d}{d\rho}\varphi(\abs{x}/\rho)\, dx
= \int \abs{\nabla v(x)}^2 \frac{\abs{x}}{(1-\theta)\rho^2}\, \chi_{A(\theta\rho, \rho)}(x)\, dx,
\]
using $(d/d\rho)\varphi(\abs{x}/\rho) =
(\abs{x}/(1-\theta)\rho^2) \chi_{[\theta\rho, \rho]}(\abs{x})$
from the piecewise-linear definition. This gives the first
equality. For the inequality, bound
$\abs{x} \geq \theta\rho$ on the annulus:
\[
\EE'(\rho) \geq \frac{\theta}{(1-\theta)\rho} \int_{A(\theta\rho, \rho)} \abs{\nabla v}^2\, dx
= \frac{\theta}{(1-\theta)\rho}\, \DD(\rho).
\]
\end{proof}

\begin{remark}[Three purposes of the hat function]\label{rem:hat}
The piecewise-linear hat function serves three distinct roles.
(i) It makes $\EE(\rho)$ monotone in $\rho$, so any bound on
$\EE$ at one scale transfers to smaller scales.
(ii) Its derivative~\eqref{eq:hat-deriv} directly controls the
annular gradient integral, giving the conversion between the
integrated energy and a pointwise-in-$\rho$ ODE quantity. (iii)
Its gradient is supported on the fixed-ratio annulus
$A(\theta\rho, \rho)$, ensuring the Bogovski\u{\i} constants
introduced in \S\ref{sec:pressure} are uniform in $\rho$.

A smooth mollification of $\varphi$ would require writing the
derivative as an average over a transition layer of width $\epsilon$;
the argument below goes through with that replacement at the cost
of an additional $\epsilon$-bookkeeping. We keep the sharp form
for clarity.
\end{remark}

\section{Hat-function energy and the iteration target}\label{sec:hat}

The proof of Theorem~\ref{thm:main} is a scale-by-scale argument
on $\EE(\rho)$. The goal (reached in \S\ref{sec:contradiction})
is the differential inequality
\begin{equation}\label{eq:target-ODE}
\EE'(\rho) \geq c\, \rho^{-\alpha}\, \EE(\rho)^\beta
\qquad \text{with } \beta > 1 \text{ and } \alpha < 1,
\end{equation}
for $\rho$ above a fixed threshold $\rho_0(M, \nu)$. Integrating
against the boundedness of $\EE$ at infinity forces
$\int_{\RR^3} \abs{\nabla v}^2\, dx < \infty$, from which a
Galdi-style testing argument (\S\ref{sec:vzero}) gives $v \equiv 0$.

The path from~\eqref{eq:hat-deriv} to~\eqref{eq:target-ODE}
consists of: pressure elimination via Bogovski\u{\i}
(\S\ref{sec:pressure}), a Caccioppoli-type inequality on annuli
(\S\ref{sec:caccioppoli}), Lorentz layer-cake embeddings
(\S\ref{sec:lorentz}), and a rescaled Gagliardo--Nirenberg
interpolation (\S\ref{sec:gn}). Each section produces a
sublinear-in-$\DD$ bound on a piece of $\EE$; combined, they give
the key estimate
\[
\EE(\rho) \leq C_0 M \rho^{-1/2} \DD(\rho)^{1/2}
          + C_1 M^3 \rho^{-1}
          + C_2 M^a \rho^\gamma \DD(\rho)^b,
\]
with explicit exponents $a = 3p/(6-p)$, $b = 3(6-2p)/(2(6-p))$,
$\gamma = (3-2p)/(6-p)$ for every $p \in (3/2, 3)$. At $p = 2$
these give $a = 6/5$, $b = 3/4$, $\gamma = -1/4$.

\section{Pressure elimination via Bogovski\u{\i}}\label{sec:pressure}

The stationary Navier--Stokes system carries a pressure field $p$
which, when tested against a non-divergence-free function,
produces a boundary contribution that does not close under
scaling. The Bogovski\u{\i} operator converts any compactly
supported cutoff of $v$ into a divergence-free competitor at the
cost of an explicit $L^q$ correction on the same support. The
pressure integral then vanishes identically, at no further cost.

\subsection{The Bogovski\u{\i} operator on John domains.}
\label{ssec:bogo}

For an open set $\Omega \subset \RR^3$ we write
$L^q_0(\Omega) := \{f \in L^q(\Omega) : \int_\Omega f\,dx = 0\}$.

\begin{lemma}[Bogovski\u{\i}]\label{lem:bogo}
Let $\Omega \subset \RR^3$ be a bounded John domain with John
constant $c_J$. There is a linear operator
\[
\BB_\Omega \colon L^q_0(\Omega) \to W^{1,q}_0(\Omega; \RR^3),
\qquad 1 < q < \infty,
\]
such that, for every $f \in L^q_0(\Omega)$,
\begin{enumerate}[label=(B\arabic*), itemsep=1pt]
\item $\dv \BB_\Omega[f] = f$ in $\Omega$, and
      $\BB_\Omega[f] \in W^{1,q}_0(\Omega; \RR^3)$;
\item $\norm{\BB_\Omega[f]}_{L^q(\Omega)} +
       \norm{\nabla \BB_\Omega[f]}_{L^q(\Omega)}
       \leq C_{q, c_J}\,\norm{f}_{L^q(\Omega)}$;
\item if $f \in L^{q_1}_0(\Omega) \cap L^{q_2}_0(\Omega)$ with
      $1 < q_1, q_2 < \infty$, then the two definitions agree;
\item $\BB_\Omega[f]$ depends linearly and continuously on $f$
      with respect to every admissible $q$.
\end{enumerate}
\end{lemma}

The construction and the bounds are due to
Bogovski\u{\i}~\cite{Bogovskii} on star-shaped domains; the
extension to John domains (in particular, to finite unions of
star-shaped sets and to annuli) is standard, see
Galdi~\cite[\S III.3]{Galdi}.

Fixed-ratio annuli are uniformly John.

\begin{lemma}[Uniform John constant for annuli]\label{lem:annulus-john}
For every $\theta \in (0, 1)$ there is a constant
$c_J(\theta)$ such that, for every $\rho > 0$, the annulus
$A(\theta\rho, \rho)$ is John with constant $c_J(\theta)$.
\end{lemma}

This is elementary: $A(\theta\rho, \rho)$ is the image of
$A(\theta, 1)$ under the dilation $y \mapsto \rho y$, and the
John condition is scale-invariant. We will use this in the
following rescaled form.

\begin{proposition}[Scaling of $\BB$ on annuli]
\label{prop:bogo-scale}
Let $\Omega_1 = A(\theta, 1)$ and $\Omega_\rho = A(\theta\rho,
\rho)$, let $f \in L^q_0(\Omega_\rho)$, and let $\tilde f(y) :=
f(\rho y)$. Then $\tilde f \in L^q_0(\Omega_1)$, and
\[
\BB_{\Omega_\rho}[f](x) = \rho\,\BB_{\Omega_1}[\tilde f](x/\rho),
\qquad
\nabla \BB_{\Omega_\rho}[f](x) = (\nabla_y \BB_{\Omega_1}[\tilde f])(x/\rho).
\]
In particular,
\begin{equation}\label{eq:bogo-scale}
\norm{\BB_{\Omega_\rho}[f]}_{L^q(\Omega_\rho)} = \rho^{1 + 3/q}
\norm{\BB_{\Omega_1}[\tilde f]}_{L^q(\Omega_1)},
\quad
\norm{\nabla \BB_{\Omega_\rho}[f]}_{L^q(\Omega_\rho)} = \rho^{3/q}
\norm{\nabla_y \BB_{\Omega_1}[\tilde f]}_{L^q(\Omega_1)}.
\end{equation}
\end{proposition}

\begin{proof}
The change of variable $y = x/\rho$ in Bogovski\u{\i}'s explicit
singular-integral representation (which is a convolution with a
kernel of degree $-(n-1)$ in $\RR^n$) shows that $\rho
\BB_{\Omega_1}[\tilde f](\cdot/\rho)$ satisfies the defining
properties (B1)--(B4) on $\Omega_\rho$ with the prescribed
operator norm. The uniqueness of the Bogovski\u{\i} operator
modulo test-function density~\cite[III.3.1]{Galdi} gives the
claimed identity. The norm formulas follow by substitution:
$\norm{\BB_{\Omega_\rho}[f]}_{L^q(\Omega_\rho)}^q =
\int_{\Omega_\rho} \rho^q \abs{\BB_{\Omega_1}[\tilde f](x/\rho)}^q
dx = \rho^{q + 3} \int_{\Omega_1}
\abs{\BB_{\Omega_1}[\tilde f](y)}^q dy$.
\end{proof}

Combining Lemmas~\ref{lem:bogo}--\ref{lem:annulus-john} and
Proposition~\ref{prop:bogo-scale} gives the uniform-in-$\rho$
norm bound that will be used throughout.

\begin{corollary}[Uniform annular Bogovski\u{\i} bound]
\label{cor:bogo-annulus}
For every $\theta \in (0, 1)$ and every $q \in (1, \infty)$ there
is a constant $C_{q,\theta}$ such that, for every $\rho > 0$ and
every $f \in L^q_0(A(\theta\rho, \rho))$,
\begin{equation}\label{eq:bogo-uniform}
\norm{\BB_{A(\theta\rho, \rho)}[f]}_{L^q} \leq C_{q,\theta}\,\rho\,\norm{f}_{L^q},
\qquad
\norm{\nabla \BB_{A(\theta\rho, \rho)}[f]}_{L^q} \leq C_{q,\theta}\,\norm{f}_{L^q}.
\end{equation}
\end{corollary}

\begin{proof}
From~\eqref{eq:bogo-scale},
$\norm{\BB_{\Omega_\rho}[f]}_{L^q} = \rho^{1+3/q}
\norm{\BB_{\Omega_1}[\tilde f]}_{L^q} \leq C_{q,\theta} \rho^{1+3/q}
\norm{\tilde f}_{L^q} = C_{q,\theta}\rho \cdot \rho^{3/q}
\norm{\tilde f}_{L^q(\Omega_1)} = C_{q,\theta}\rho \norm{f}_{L^q}$,
using $\norm{f}^q_{L^q(\Omega_\rho)} = \rho^3 \norm{\tilde f}^q_{L^q(\Omega_1)}$.
Similarly for $\nabla \BB$.
\end{proof}

\subsection{The divergence-free test function.}

Fix $\theta \in (0, 1)$ throughout. For $\rho > 0$, let
$\eta \in C^\infty_c(\RR^3)$ be a radial cutoff function with
\begin{equation}\label{eq:cutoff}
\eta \equiv 1 \text{ on } B_{\theta\rho}, \quad
\supp \eta \subset B_\rho, \quad
0 \leq \eta \leq 1, \quad
\abs{\nabla \eta} \leq \frac{C_\theta}{\rho}, \quad
\abs{\Delta \eta} \leq \frac{C_\theta}{\rho^2},
\end{equation}
and $\supp \nabla \eta \subset A(\theta\rho, \rho)$. The constant
$C_\theta$ depends only on $\theta$. (Such an $\eta$ is standard:
take a smooth mollification of $\varphi(\cdot/\rho)^2$ supported in
the annulus, or any radial bump satisfying the displayed bounds.)

Define the auxiliary field
\begin{equation}\label{eq:bogo-aux}
b(x) := \BB_{A(\theta\rho, \rho)}\!\bigl[v(\cdot) \cdot \nabla\eta(\cdot)\bigr](x),
\qquad x \in \RR^3,
\end{equation}
extended by zero outside $A(\theta\rho, \rho)$. The argument
$v \cdot \nabla \eta$ is integrable and has mean zero on
$A(\theta\rho, \rho)$ because
\[
\int_{A(\theta\rho,\rho)} v \cdot \nabla \eta\, dx =
\int_{\RR^3} v \cdot \nabla \eta\, dx =
-\int_{\RR^3} \eta\, (\dv v)\, dx = 0,
\]
so $b$ is well defined by Lemma~\ref{lem:bogo}.

\begin{lemma}[Test-function properties]\label{lem:w-props}
The field $w := v\,\eta - b$ is divergence-free on $\RR^3$,
lies in $W^{1,q}(\RR^3; \RR^3)$ for every $q \in (1, \infty)$,
and is compactly supported in $\overline{B_\rho}$.
\end{lemma}

\begin{proof}
Support: $v\eta$ is supported in $\supp \eta \subset B_\rho$, and
$b$ is supported in $A(\theta\rho, \rho) \subset B_\rho$ by (B1).

Divergence:
\[
\dv(v\eta) = \eta\,\dv v + v \cdot \nabla \eta = v \cdot \nabla \eta,
\]
and by (B1), $\dv b = v \cdot \nabla \eta$ on
$A(\theta\rho, \rho)$; since both functions are zero outside
$A(\theta\rho, \rho)$, we have $\dv b = v \cdot \nabla \eta$ on
all of $\RR^3$ in the sense of distributions. Subtracting,
$\dv w = 0$.

Regularity: $v$ and $\eta$ are smooth, hence $v\eta \in C^\infty_c$,
and $b \in W^{1,q}_0(A(\theta\rho, \rho); \RR^3)$ by (B2) for
every $q$.
\end{proof}

\subsection{The resulting identity.}

\begin{proposition}[Pressure-eliminated identity]\label{prop:elim}
With $\eta$ and $w$ as above,
\begin{equation}\label{eq:master-identity}
\nu\int_{\RR^3} \abs{\nabla v}^2\,\eta\, dx
= \frac{\nu}{2}\int_{\RR^3} \abs{v}^2 \Delta\eta\, dx
+ \frac{1}{2}\int_{\RR^3} \abs{v}^2\,(v \cdot \nabla \eta)\, dx
- \int_{\RR^3} (v \otimes v) \colon \nabla b\, dx
+ \nu\int_{\RR^3} \nabla v \colon \nabla b\, dx.
\end{equation}
\end{proposition}

\begin{proof}
Pair the stationary Navier--Stokes equation
$-\nu\Delta v + (v \cdot \nabla) v + \nabla p = 0$ with
$w = v\eta - b$ and integrate over $\RR^3$. Since $\dv w = 0$ and
$w$ is compactly supported,
\[
\int \nabla p \cdot w\, dx = -\int p\, \dv w\, dx = 0.
\]
So the identity reduces to
\begin{equation}\label{eq:pair-basic}
-\nu \int (\Delta v) \cdot w\, dx + \int (v \cdot \nabla v) \cdot w\, dx = 0.
\end{equation}
Expand $w = v\eta - b$ and integrate by parts (valid by the
compact support of $w$ and smoothness of $v$):
\[
-\int (\Delta v)_i \cdot (v_i \eta)\, dx
= \int (\partial_j v_i)\,\partial_j(v_i \eta)\, dx
= \int \abs{\nabla v}^2 \eta\, dx
+ \frac{1}{2}\int \nabla \abs{v}^2 \cdot \nabla \eta\, dx,
\]
and a final integration by parts on the last term gives
$-\frac{1}{2}\int \abs{v}^2 \Delta \eta\, dx$. Hence
\[
-\nu \int (\Delta v) \cdot (v\eta)\, dx
= \nu \int \abs{\nabla v}^2 \eta\, dx - \frac{\nu}{2}\int \abs{v}^2 \Delta \eta\, dx.
\]
For the $b$-piece,
\[
-\nu \int (\Delta v) \cdot (-b)\, dx = -\nu \int \nabla v \colon \nabla b\, dx,
\]
by integration by parts (using $b \in W^{1,q}_0$ for every $q$).
For the convective pairing against $v\eta$, use $\dv v = 0$:
\[
\int (v \cdot \nabla v)_i (v_i \eta)\, dx
= \frac{1}{2}\int v \cdot \nabla \abs{v}^2 \, \eta \, dx
= -\frac{1}{2}\int \abs{v}^2\,(v \cdot \nabla \eta)\, dx,
\]
the final IBP using $\dv v = 0$ again. For the convective pairing
against $-b$, use $\dv v = 0$ once more:
\[
-\int (v \cdot \nabla v)_i\, b_i\, dx
= \int v_i v_j\, \partial_j b_i\, dx
= \int (v \otimes v) \colon \nabla b\, dx.
\]
Substituting the four pieces into~\eqref{eq:pair-basic} gives
\[
\Bigl[\nu \int \abs{\nabla v}^2 \eta - \frac{\nu}{2}\int \abs{v}^2 \Delta \eta\Bigr]
- \nu \int \nabla v \colon \nabla b
+ \Bigl[-\frac{1}{2}\int \abs{v}^2 v \cdot \nabla \eta\Bigr]
+ \int (v\otimes v) \colon \nabla b = 0,
\]
and rearrangement yields~\eqref{eq:master-identity}.
\end{proof}

\begin{remark}[Why the Bogovski\u{\i} correction is forced]
\label{rem:why-bogo}
A naive test function $v\eta$ is not divergence-free unless
$\eta$ is constant on $\supp v$; the correction
$\BB[v\cdot\nabla\eta]$ restores the constraint exactly. Without
it, $\int \nabla p \cdot v\eta\, dx = -\int p\,v\cdot\nabla\eta\,
dx$ carries a pressure remainder that cannot be bounded uniformly
in $\rho$ by any standard argument at the critical exponent. The
trade is paid in the two new terms
$\nu \int \nabla v \colon \nabla b$ and $\int (v\otimes v) \colon
\nabla b$ on the right of~\eqref{eq:master-identity}; both are
controlled in the next section.
\end{remark}

\section{Caccioppoli inequality on annuli}\label{sec:caccioppoli}

From the identity~\eqref{eq:master-identity} we extract a
Caccioppoli-type bound whose right-hand side is \emph{sublinear}
in $\DD(\rho)$. The restructuring avoids the Young-inequality
absorption step that, in earlier drafts, forced a case
dichotomy; it keeps cross terms as products
$\norm{\nabla v}_{L^2} \cdot (\text{scale-corrected factor})$
bounded by direct Cauchy--Schwarz plus Lorentz embedding.

Throughout this section we take
$\eta = \varphi(\abs{\cdot}/\rho)$ as the hat function of
Notation~\ref{not:geometry}. The identity
of Proposition~\ref{prop:elim} holds for this choice as the
limit, as $\varepsilon \to 0$, of the smooth-mollification
identity for $\eta_\varepsilon \in C^\infty_c$ satisfying
\eqref{eq:cutoff}; both sides of~\eqref{eq:master-identity}
converge by dominated convergence and the Bogovski\u{\i}
linearity~\cite[III.3.1]{Galdi}. With this choice,
\[
\int_{\RR^3} \abs{\nabla v}^2 \eta\, dx = \EE(\rho).
\]

\begin{proposition}[Sublinear Caccioppoli estimate]
\label{prop:caccioppoli}
There is a constant $C_\star = C_\star(\theta, \nu, M)$ such that
for every $\rho > 0$,
\begin{equation}\label{eq:cacc-full}
\nu \EE(\rho) \leq C_\star \rho^{-1/2} \DD(\rho)^{1/2}
+ \frac{\nu}{2}\abs{I_1(\rho)}
+ \frac{1}{2}\int_{\RR^3} \abs{v}^2 \abs{v\cdot\nabla\eta}\, dx,
\end{equation}
where the residual viscous term
$\abs{I_1(\rho)} \leq C_\theta \nu M^2 / \rho$, and the third
term (the convective cutoff) is bounded by
$C_\theta \rho^{-1} \norm{v}_{L^3(A(\theta\rho, \rho))}^3$.
\end{proposition}

\begin{proof}
By~\eqref{eq:master-identity} (rearranged):
\[
\nu \EE(\rho) \leq \frac{\nu}{2}\abs{\int \abs{v}^2 \Delta\eta\, dx}
+ \frac{1}{2}\abs{\int\abs{v}^2 v \cdot \nabla\eta\, dx}
+ \abs{\int (v\otimes v) : \nabla b\, dx}
+ \nu \abs{\int \nabla v : \nabla b\, dx}.
\]

We bound each of the four terms.

\emph{Term $I_1$ (viscous cutoff).} Using $\abs{\Delta\eta} \leq
C_\theta/\rho^2$ on $A(\theta\rho, \rho)$ (supp of $\Delta\eta$),
\[
\abs{I_1} := \abs{\int \abs{v}^2 \Delta\eta\, dx}
\leq \frac{C_\theta}{\rho^2} \norm{v}_{L^2(A(\theta\rho, \rho))}^2.
\]
By the Lorentz layer-cake embedding proven in~\S\ref{sec:lorentz}
(Proposition~\ref{prop:lorentz-emb}), $\norm{v}_{L^2(A)}^2 \leq
C\,M^2\,\rho$, so
$\abs{I_1} \leq C_\theta M^2 / \rho$, and
$\frac{\nu}{2}\abs{I_1} \leq C_\theta \nu M^2/\rho$.

\emph{Term $I_2$ (Bogovski\u{\i} viscous).} By Cauchy--Schwarz,
\[
\abs{I_2} := \abs{\int \nabla v : \nabla b\, dx}
\leq \norm{\nabla v}_{L^2(A)} \cdot \norm{\nabla b}_{L^2(A)}
= \DD(\rho)^{1/2} \cdot \norm{\nabla b}_{L^2(A)}.
\]
Using the uniform annular Bogovski\u{\i} bound
Corollary~\ref{cor:bogo-annulus} with $q = 2$ and
$\abs{\nabla \eta} \leq C_\theta/\rho$,
\[
\norm{\nabla b}_{L^2(A)} \leq C_\theta \norm{v \cdot \nabla \eta}_{L^2(A)}
\leq \frac{C_\theta}{\rho} \norm{v}_{L^2(A)}
\leq \frac{C_\theta M}{\sqrt{\rho}},
\]
by the Lorentz bound of~\S\ref{sec:lorentz}. Combining,
\[
\nu \abs{I_2} \leq C_\theta \nu M\, \rho^{-1/2}\,\DD(\rho)^{1/2}.
\]

\emph{Term $I_4$ (convective Bogovski\u{\i}).} Apply the
three-factor H\"older inequality with exponents $(12/5, 2, 12)$
to the identity $\int (v\otimes v):\nabla b\, dx =
\int v_i\,v_j\,\partial_j b_i\, dx$; equivalently (via the
integration-by-parts identity $\int (v\otimes v):\nabla b =
-\int (v \cdot \nabla v)\cdot b$ from
Proposition~\ref{prop:elim}), bound
\[
\abs{I_4} := \abs{\int (v \otimes v) : \nabla b\, dx}
= \abs{\int (v \cdot \nabla v) \cdot b\, dx}
\leq \norm{v}_{L^{12/5}(A)}\,\norm{\nabla v}_{L^2(A)}\,\norm{b}_{L^{12}(A)}.
\]
By the Sobolev embedding
$W^{1,12/5}_0(A) \hookrightarrow L^{12}(A)$ (rescaled from the
reference annulus $A_0 = A(\theta, 1)$, with constant depending
only on $\theta$),
\[
\norm{b}_{L^{12}(A)} \leq C_\theta \norm{\nabla b}_{L^{12/5}(A)}.
\]
Bogovski\u{\i} at $q = 12/5$ (Corollary~\ref{cor:bogo-annulus})
gives
\[
\norm{\nabla b}_{L^{12/5}(A)} \leq C_\theta \norm{v \cdot \nabla \eta}_{L^{12/5}(A)}
\leq \frac{C_\theta}{\rho} \norm{v}_{L^{12/5}(A)}.
\]
By the Lorentz embedding (Proposition~\ref{prop:lorentz-emb})
at $r = 12/5$,
\[
\norm{v}_{L^{12/5}(A)} \leq C M \rho^{3(5/12) - 1} = C M \rho^{1/4}.
\]
Assembling,
\[
\abs{I_4} \leq (C M \rho^{1/4}) \cdot \DD(\rho)^{1/2} \cdot
\Bigl(\frac{C_\theta}{\rho} \cdot C M \rho^{1/4}\Bigr)
= C_\theta M^2 \rho^{-1/2} \DD(\rho)^{1/2}.
\]

Setting $C_\star := C_\theta(\nu M + M^2)$ (so that
$\nu\abs{I_2} + \abs{I_4} \leq C_\star \rho^{-1/2} \DD^{1/2}$)
and collecting the bounds on $I_1$, $I_2$, $I_4$, we
obtain~\eqref{eq:cacc-full}. The residual $I_3$-type term
$\frac{1}{2}\int \abs{v}^2 v \cdot \nabla\eta$ is bounded above
by $\frac{C_\theta}{\rho}\norm{v}_{L^3(A)}^3$ using
$\abs{\nabla\eta} \leq C_\theta/\rho$ on the support $A$.
\end{proof}

The third term of~\eqref{eq:cacc-full} involves the scale-critical
quantity $\norm{v}_{L^3(A)}^3$, which is not bounded directly by
Lorentz (the Lorentz embedding of~\S\ref{sec:lorentz} requires
$r < 3$ strictly). This term is handled by the Gagliardo--Nirenberg
interpolation of~\S\ref{sec:gn}, which factors it through the
sub-scale-critical Lorentz bounds plus the gradient norm.

\begin{remark}[Draft-9 sublinearity vs.\ Young absorption]
\label{rem:no-young}
Earlier approaches to stationary Liouville (Galdi's original
argument, and Drafts~7--8 in the companion corpus) attempted to
absorb the Bogovski\u{\i}-viscous cross-term into the left-hand
side by Young's inequality: $ab \leq \tfrac{1}{2}a^2 +
\tfrac{1}{2}b^2$ with the $a^2$-piece absorbed into $\DD(\rho)$.
This requires $\DD$ to appear on the left of~\eqref{eq:cacc-full},
which is not the case here: the left-hand side is $\EE(\rho)$,
not $\DD$. The identity~\eqref{eq:master-identity} cannot be
rearranged to put $\DD$ on the left without introducing a
Caccioppoli inequality for $\DD$ itself, which is not available
at this regularity. The sublinear form~\eqref{eq:cacc-full} with
$\DD^{1/2}$ on the right is therefore the correct structural
form, confirmed by the explicit Bogovski\u{\i} scaling. The
retired Young-absorption approach is catalogued as anti-knowledge
AK-28 in the companion corpus.
\end{remark}

\section{Lorentz layer-cake embedding on annuli}\label{sec:lorentz}

Recall $\norm{f}_{L^{q,\infty}(\RR^n)} := \sup_{t > 0} t\,
\abs{\{\abs{f} > t\}}^{1/q}$. For $f \in L^{q,\infty}$ restricted
to a set $E$ of finite measure, interpolation with $L^\infty$
gives an $L^r$ bound for every $r \in (0, q)$.

\begin{proposition}[Lorentz embedding on sets of finite measure]
\label{prop:lorentz-emb}
Let $f \in L^{q,\infty}(\RR^n)$ with $q \in (0, \infty)$, and let
$E \subset \RR^n$ be a measurable set of finite measure. For
every $r \in (0, q)$,
\begin{equation}\label{eq:lorentz}
\norm{f}_{L^r(E)} \leq \Bigl(\frac{q}{q - r}\Bigr)^{1/r}
\norm{f}_{L^{q,\infty}(\RR^n)}\, \abs{E}^{1/r - 1/q}.
\end{equation}
\end{proposition}

\begin{proof}
Write $M := \norm{f}_{L^{q,\infty}}$. By Cavalieri's identity,
\[
\int_E \abs{f}^r\, dx = r \int_0^\infty t^{r-1}\,
\abs{\{x \in E : \abs{f(x)} > t\}}\, dt.
\]
Bound the distribution function by the minimum of its two
available estimates:
\[
\abs{\{x \in E : \abs{f(x)} > t\}} \leq \min\Bigl(\abs{E},\, M^q / t^q\Bigr).
\]
The two bounds are equal at $t = t_* := (M^q / \abs{E})^{1/q}$.
Splitting the integral at $t_*$,
\[
\int_E \abs{f}^r\, dx \leq r\,\abs{E}\int_0^{t_*} t^{r-1}\,dt
+ r\,M^q \int_{t_*}^\infty t^{r-1-q}\, dt
= \abs{E}\, t_*^r + \frac{r}{q-r}\, M^q\, t_*^{r - q}.
\]
Substituting $t_*^r = (M^q/\abs{E})^{r/q}$ and $t_*^{r - q} =
(\abs{E}/M^q)^{1 - r/q}$,
\[
\int_E \abs{f}^r\, dx \leq M^r \abs{E}^{1 - r/q}
\Bigl(1 + \frac{r}{q - r}\Bigr)
= \frac{q}{q - r}\, M^r\, \abs{E}^{1 - r/q}.
\]
Taking the $r$-th root yields~\eqref{eq:lorentz}.
\end{proof}

We apply~\eqref{eq:lorentz} with $q = 3$, $f = v$, and $E =
A(\theta\rho, \rho)$. The volume of the annulus is
\[
\abs{A(\theta\rho, \rho)} = \tfrac{4\pi}{3}(1 - \theta^3)\rho^3
=: c_\theta \rho^3.
\]

\begin{corollary}[Lorentz bound on the annulus]
\label{cor:lorentz-v}
For every $r \in (0, 3)$ and every $\rho > 0$,
\begin{equation}\label{eq:v-lorentz}
\norm{v}_{L^r(A(\theta\rho, \rho))}
\leq C_{r, \theta}\, M\, \rho^{3/r - 1}.
\end{equation}
\end{corollary}

\begin{proof}
Apply Proposition~\ref{prop:lorentz-emb} with $q = 3$:
$\norm{v}_{L^r(A)} \leq (3/(3-r))^{1/r} M \cdot (c_\theta
\rho^3)^{1/r - 1/3} = C_{r,\theta} M \rho^{3/r - 1}$.
\end{proof}

\begin{remark}[Why the embedding fails at $r = 3$]
\label{rem:lorentz-endpoint}
As $r \to 3^-$, the prefactor $(3/(3-r))^{1/r}$ diverges, and
the exponent $3/r - 1 \to 0$. The Lorentz quasi-norm
$\norm{\cdot}_{L^{3,\infty}}$ is not absolutely continuous:
a function like $v(x) = \abs{x}^{-1}$ has
$\norm{v}_{L^{3,\infty}(\RR^3)} < \infty$ but
$\norm{v}_{L^3(B_R \setminus B_r)} = c\log^{1/3}(R/r) \to \infty$
as $r \to 0$. This is the structural reason Theorem~\ref{thm:main}
requires the sub-scale-critical machinery of~\S\ref{sec:gn}: at
$r = 3$, no Lorentz-based bound is available, and the critical
exponent must be accessed through the gradient via
Gagliardo--Nirenberg.
\end{remark}

The key bounds used in~\S\ref{sec:caccioppoli} are the specific
instances $r = 2$ (giving $\norm{v}_{L^2(A)} \leq CM\rho^{1/2}$)
and $r = 12/5$ (giving $\norm{v}_{L^{12/5}(A)} \leq
CM\rho^{1/4}$). The sub-scale-critical range $r \in (3/2, 3)$
underlies~\S\ref{sec:gn}.

\section{Gagliardo--Nirenberg interpolation and scale cancellation}
\label{sec:gn}

The convective cutoff term in~\eqref{eq:cacc-full} is cubic in
$v$ and requires an $L^3(A)$ bound on $v$. The Lorentz embedding
of~\S\ref{sec:lorentz} does not reach $r = 3$ endpoint. We
therefore interpolate $\norm{v}_{L^3(A)}^3$ between a
sub-critical Lorentz bound and the annular gradient norm, via
Gagliardo--Nirenberg, and observe that the exact scale-invariance
of the critical $L^3$ exponent forces a cancellation of
$\rho$-powers in the main term.

\subsection{Reference Gagliardo--Nirenberg inequality.}

\begin{lemma}[Gagliardo--Nirenberg on the reference annulus]
\label{lem:GN-ref}
Fix $\theta \in (0, 1)$ and $p \in (3/2, 3)$, and set $\theta_p :=
p/(6 - p) \in (0, 1)$. There is a constant $C_{\mathrm{GN}}(p, \theta)$ such
that for every $w \in W^{1,2}(A(\theta, 1)) \cap L^p(A(\theta, 1))$,
\begin{equation}\label{eq:GN-ref}
\norm{w}_{L^3(A(\theta,1))} \leq C_{\mathrm{GN}}\bigl(
\norm{w}_{L^p(A(\theta,1))}^{\theta_p}
\norm{\nabla w}_{L^2(A(\theta,1))}^{1 - \theta_p}
+ \norm{w}_{L^p(A(\theta,1))} \bigr).
\end{equation}
\end{lemma}

\begin{proof}
This is the standard multiplicative Gagliardo--Nirenberg
inequality on a bounded Lipschitz domain~\cite[Thm~II.5.2]{Galdi}.
The exponent $\theta_p$ is determined by the scaling relation
\[
\frac{1}{3} = \frac{\theta_p}{p} + (1 - \theta_p)\Bigl(\frac{1}{2} - \frac{1}{3}\Bigr),
\]
which rearranges to $\theta_p = p/(6 - p)$. The additive
$\norm{w}_{L^p}$ term is required on bounded domains because
$\norm{w}_{L^p} + \norm{\nabla w}_{L^2}$ bounds the full
$W^{1,p} \cap W^{1,2}$ norm only up to this lower-order piece.
The constant $C_{\mathrm{GN}}$ depends on the annulus $A(\theta, 1)$ through
standard Sobolev-extension estimates.
\end{proof}

Cubing~\eqref{eq:GN-ref} and using $(x + y)^3 \leq 4(x^3 + y^3)$,
\begin{equation}\label{eq:GN-cubed}
\norm{w}_{L^3(A(\theta,1))}^3 \leq C'_{\mathrm{GN}}\Bigl(
\norm{w}_{L^p(A(\theta,1))}^a\,\norm{\nabla w}_{L^2(A(\theta,1))}^{2b}
+ \norm{w}_{L^p(A(\theta,1))}^3 \Bigr),
\end{equation}
where
\[
a := 3\theta_p = \frac{3p}{6 - p}, \qquad
b := \frac{3(1 - \theta_p)}{2} = \frac{3(3 - p)}{6 - p}, \qquad
a + 2b = 3.
\]

\subsection{Rescaling to $A(\theta\rho, \rho)$ and scale cancellation.}

\begin{proposition}[Critical GN bound with exact $\rho$-cancellation]
\label{prop:GN-critical}
For every $\theta \in (0, 1)$ and $p \in (3/2, 3)$, and every
$\rho > 0$,
\begin{equation}\label{eq:GN-scaled}
\norm{v}_{L^3(A(\theta\rho, \rho))}^3
\leq C'_{\mathrm{GN}} \Bigl(
\norm{v}_{L^p(A(\theta\rho, \rho))}^a\, \DD(\rho)^b
+ \norm{v}_{L^p(A(\theta\rho, \rho))}^3 \Bigr),
\end{equation}
with no explicit $\rho$-factor. Substituting the Lorentz bound
Corollary~\ref{cor:lorentz-v} yields
\begin{equation}\label{eq:GN-Lorentz}
\norm{v}_{L^3(A(\theta\rho, \rho))}^3 \leq
C_{\mathrm{GN}}'(CM)^a\,\rho^b\, \DD(\rho)^b + C_{\mathrm{GN}}'(CM)^3.
\end{equation}
\end{proposition}

\begin{proof}
Let $\tilde v(y) := v(\rho y)$ for $y \in A(\theta, 1)$. Direct
substitution in each integral gives
\[
\norm{\tilde v}_{L^q(A(\theta,1))}^q = \rho^{-3}\,\norm{v}_{L^q(A(\theta\rho,\rho))}^q,
\qquad
\norm{\nabla_y \tilde v}_{L^2(A(\theta,1))}^2 = \rho^{-1}\,\DD(\rho),
\]
for any $q \geq 1$. Substituting into~\eqref{eq:GN-cubed} applied
to $\tilde v$,
\[
\rho^{-3}\,\norm{v}_{L^3(A)}^3
\leq C'_{\mathrm{GN}}\bigl(
\rho^{-3a/p - b}\norm{v}_{L^p(A)}^a\,\DD(\rho)^b
+ \rho^{-9/p}\norm{v}_{L^p(A)}^3
\bigr).
\]
Multiplying by $\rho^3$, the first prefactor becomes
$\rho^{3 - 3a/p - b}$. Substitute $a = 3p/(6 - p)$, $b = 3(3 - p)/(6 - p)$:
\[
3 - \frac{3a}{p} - b
= 3 - \frac{9}{6-p} - \frac{3(3 - p)}{6 - p}
= \frac{3(6-p) - 9 - 9 + 3p}{6 - p}
= \frac{18 - 3p - 18 + 3p}{6 - p} = 0.
\]
The second prefactor becomes
\[
\rho^{3 - 9/p}\,\norm{v}_{L^p(A)}^3 \leq
\rho^{3 - 9/p}\cdot(CM)^3\rho^{3(3/p - 1)}
= (CM)^3\,\rho^{3 - 9/p + 9/p - 3} = (CM)^3,
\]
using Lorentz $\norm{v}_{L^p(A)} \leq CM\rho^{3/p - 1}$. This
proves~\eqref{eq:GN-scaled} and, with the Lorentz substitution
$\norm{v}_{L^p(A)}^a \leq (CM)^a \rho^{a(3/p - 1)}$ where
$a(3/p - 1) = b$ by direct computation, yields
\eqref{eq:GN-Lorentz}.
\end{proof}

\subsection{The full key estimate.}

\begin{theorem}[Main sublinear estimate on $\EE$]
\label{thm:main-estimate}
Fix $\theta \in (0, 1)$ and $p \in (3/2, 3)$, with $\gamma :=
(3 - 2p)/(6 - p)$ and $b := 3(3 - p)/(6 - p)$. There are
constants $K_0, K_1, K_2$ depending only on $(\theta, \nu, M, p)$
such that, for every $\rho > 0$,
\begin{equation}\label{eq:key}
\EE(\rho) \leq K_0\,\rho^{-1/2}\,\DD(\rho)^{1/2}
          + K_1\,\rho^{-1}
          + K_2\,\rho^\gamma\,\DD(\rho)^b.
\end{equation}
\end{theorem}

\begin{proof}
Divide~\eqref{eq:cacc-full} (Proposition~\ref{prop:caccioppoli})
by $\nu$:
\[
\EE(\rho) \leq \frac{C_\star}{\nu}\,\rho^{-1/2}\DD^{1/2}
+ \frac{1}{2}\abs{I_1(\rho)}
+ \frac{1}{2\nu}\int \abs{v}^2\abs{v\cdot\nabla\eta}\,dx.
\]
The first term contributes $K_0\rho^{-1/2}\DD^{1/2}$ with
$K_0 = C_\star/\nu$.

The second term is bounded by
$\abs{I_1} \leq C_\theta M^2/\rho$, contributing
$C_\theta M^2/(2\rho)$.

The third term, using $\abs{\nabla\eta} \leq C_\theta/\rho$ on
$A(\theta\rho, \rho)$,
\[
\int \abs{v}^2 \abs{v\cdot\nabla\eta}\,dx
\leq \frac{C_\theta}{\rho}\,\norm{v}_{L^3(A(\theta\rho,\rho))}^3.
\]
Applying~\eqref{eq:GN-Lorentz},
\[
\frac{1}{\rho}\,\norm{v}_{L^3(A)}^3
\leq C'_{\mathrm{GN}}(CM)^a\,\rho^{b - 1}\,\DD^b
+ \frac{C'_{\mathrm{GN}}(CM)^3}{\rho}.
\]
Since $b - 1 = \gamma$,
\[
\frac{1}{2\nu}\int \abs{v}^2\abs{v\cdot\nabla\eta}\,dx
\leq \frac{C_\theta C'_{\mathrm{GN}} C^a M^a}{2\nu}\,\rho^\gamma\,\DD^b
+ \frac{C_\theta C'_{\mathrm{GN}} C^3 M^3}{2\nu\rho}.
\]
Combining the $\rho^{-1}$ contributions into $K_1$, the
$\rho^\gamma\DD^b$ contribution into $K_2$, and the
$\rho^{-1/2}\DD^{1/2}$ contribution into $K_0$
gives~\eqref{eq:key}.
\end{proof}

\begin{remark}[Exponents at $p = 2$]\label{rem:exponents}
The natural choice $p = 2$ gives
\[
a = \frac{6}{4} = \frac{3}{2}, \quad
b = \frac{3}{4}, \quad
\gamma = -\frac{1}{4}.
\]
The exponents $(\gamma, b) = (-1/4, 3/4)$ are the ones used in
the Draft-9 closure of the Caccioppoli restructuring in the
companion corpus; they give the cleanest integration in
\S\ref{sec:ode}.
\end{remark}

\section{The superlinear differential inequality}\label{sec:ode}

We reduce Theorem~\ref{thm:main-estimate} to a closed ODE for
$\EE$. Throughout this section we fix $p \in (3/2, 12/5]$ so
that the critical GN exponent $b = 3(3-p)/(6-p) \geq 1/2$, and
we will focus on the canonical choice $p = 2$ for the explicit
exponents.

\begin{proposition}[Superlinear ODE]\label{prop:ODE}
Fix $p \in (3/2, 12/5]$, with $\gamma = (3 - 2p)/(6 - p)$ and
$b = 3(3 - p)/(6 - p)$. If $\lim_{\rho \to \infty} \EE(\rho) =
+\infty$, then there exist $c > 0$ and $\rho_* \geq 1$, depending
on $(\theta, \nu, M, p)$, such that
\begin{equation}\label{eq:superlinear-ODE}
\EE'(\rho) \geq c\,\rho^{-\alpha}\,\EE(\rho)^\beta
\qquad \text{for all } \rho \geq \rho_*,
\end{equation}
with
\[
\alpha = 1 + \gamma/b = \frac{12 - 5p}{3(3 - p)},
\qquad
\beta = \frac{1}{b} = \frac{6 - p}{3(3 - p)}.
\]
Moreover, $\alpha < 1$ and $\beta > 1$ for every
$p \in (3/2, 3)$.
\end{proposition}

\begin{proof}
Since $b \geq 1/2$, for any $\DD \geq 0$,
\[
\DD^{1/2} \leq \DD^b + 1,
\]
distinguishing $\DD \geq 1$ (where $\DD^{1/2} \leq \DD^b$) from
$\DD < 1$ (where $\DD^{1/2} \leq 1$). Substituting
into~\eqref{eq:key}:
\begin{equation}\label{eq:key2}
\EE(\rho) \leq K_0\rho^{-1/2}(\DD(\rho)^b + 1)
+ K_1\rho^{-1}
+ K_2\rho^\gamma\DD(\rho)^b
= (K_0\rho^{-1/2} + K_2\rho^\gamma)\DD(\rho)^b + R(\rho),
\end{equation}
where $R(\rho) := K_0\rho^{-1/2} + K_1\rho^{-1}$. Since
$\gamma \geq -1/2$ for $p \in (3/2, 12/5]$ (explicit check:
$\gamma = -1/2$ at $p = 12/5$; $\gamma$ decreases with $p$
from 0 at $p = 3/2$), for $\rho \geq 1$ we have
$\rho^{-1/2} \leq \rho^\gamma$, so
\[
\EE(\rho) \leq (K_0 + K_2)\rho^\gamma \DD(\rho)^b + R(\rho).
\]
Set $K := K_0 + K_2$.

By the hypothesis $\EE \to \infty$, there is $\rho_1 \geq 1$
such that $\EE(\rho) \geq 2R(\rho)$ for all $\rho \geq \rho_1$
(since $R(\rho) \to 0$ and $\EE(\rho) \to \infty$). Then
\[
\EE(\rho) - R(\rho) \geq \tfrac{1}{2}\EE(\rho),
\quad\text{so}\quad
\tfrac{1}{2}\EE(\rho) \leq K\rho^\gamma \DD(\rho)^b
\quad\forall \rho \geq \rho_1.
\]
Rearranging,
\[
\DD(\rho) \geq \Bigl(\frac{\EE(\rho)}{2K}\Bigr)^{1/b}
\rho^{-\gamma/b}.
\]
By Lemma~\ref{lem:hat-deriv},
$\DD(\rho) \leq \frac{1-\theta}{\theta}\,\rho\,\EE'(\rho)$, so
\[
\rho\,\EE'(\rho) \geq \frac{\theta}{1 - \theta}
\Bigl(\frac{\EE(\rho)}{2K}\Bigr)^{1/b} \rho^{-\gamma/b}.
\]
Divide by $\rho$:
\[
\EE'(\rho) \geq c\,\rho^{-1-\gamma/b}\,\EE(\rho)^{1/b},
\qquad c := \frac{\theta}{(1-\theta)(2K)^{1/b}}.
\]
This is~\eqref{eq:superlinear-ODE} with $\alpha = 1 + \gamma/b$
and $\beta = 1/b$, valid for $\rho \geq \rho_* := \rho_1$.

It remains to verify $\alpha < 1$ and $\beta > 1$. Substituting
$\gamma = (3-2p)/(6-p)$ and $b = 3(3-p)/(6-p)$,
\[
\frac{\gamma}{b}
= \frac{(3-2p)/(6-p)}{3(3-p)/(6-p)}
= \frac{3 - 2p}{3(3 - p)},
\qquad
\alpha = 1 + \frac{3 - 2p}{3(3-p)}
= \frac{3(3-p) + 3 - 2p}{3(3-p)}
= \frac{12 - 5p}{3(3 - p)},
\]
\[
\beta = \frac{6 - p}{3(3 - p)}.
\]
For $p \in (3/2, 3)$, the denominator $3(3 - p) > 0$. Then
\[
\alpha < 1 \iff 12 - 5p < 3(3 - p) = 9 - 3p \iff 2p > 3 \iff p > 3/2,
\]
\[
\beta > 1 \iff 6 - p > 3(3 - p) = 9 - 3p \iff 2p > 3 \iff p > 3/2.
\]
Both conditions hold exactly when $p > 3/2$.
\end{proof}

\begin{corollary}[Explicit exponents at $p = 2$]\label{cor:p2-expo}
At $p = 2$,
\[
\gamma = -\tfrac{1}{4}, \quad b = \tfrac{3}{4},
\quad \alpha = \tfrac{2}{3}, \quad \beta = \tfrac{4}{3}.
\]
\end{corollary}

\begin{remark}[Sharp failure at $p = 3/2$ and $p = 3$]
The endpoint $p = 3/2$ corresponds to $\alpha = \beta = 1$, where
the superlinear structure degenerates (the ODE becomes linear
and~\S\ref{sec:contradiction} does not close). The endpoint
$p = 3$ lies outside the Lorentz embedding range of
\S\ref{sec:lorentz}. The open range $p \in (3/2, 3)$ is therefore
the full admissible set; the argument is strongest at $p$ away
from both endpoints.
\end{remark}

\section{Integration and contradiction}\label{sec:contradiction}

We integrate the ODE of Proposition~\ref{prop:ODE} and derive a
contradiction with the non-decreasing, finite-at-each-scale
behaviour of $\EE$.

\begin{theorem}[Finite Dirichlet energy]\label{thm:finite-D}
$\int_{\RR^3} \abs{\nabla v}^2\, dx < \infty$.
\end{theorem}

\begin{proof}
By Lemma~\ref{lem:hat-deriv}, $\EE$ is non-decreasing, so either
$\lim_{\rho \to \infty}\EE(\rho) = L < \infty$ or $\lim_{\rho \to
\infty}\EE(\rho) = +\infty$. In the first case,
\[
\int_{\RR^3}\abs{\nabla v}^2\,dx
= \lim_{\rho \to \infty} \int \abs{\nabla v}^2\,\varphi(\abs{x}/\rho)\,dx
= L < \infty,
\]
by monotone convergence ($\varphi(\abs{x}/\rho) \nearrow 1$
pointwise as $\rho \to \infty$).

It remains to rule out the second case. Assume, for contradiction,
$\EE(\rho) \to +\infty$. Proposition~\ref{prop:ODE} applies: for
some $c > 0$, $\rho_* \geq 1$, and $(\alpha, \beta)$ with
$\alpha < 1 < \beta$ (at $p = 2$: $\alpha = 2/3$, $\beta = 4/3$),
\[
\EE'(\rho) \geq c\,\rho^{-\alpha}\,\EE(\rho)^\beta
\qquad \forall \rho \geq \rho_*.
\]
Since $\EE(\rho) > 0$ (on the hypothesis $\EE \to \infty$ we may
assume $\EE(\rho_*) > 0$), divide by $\EE^\beta$:
\[
\EE(\rho)^{-\beta}\,\EE'(\rho) \geq c\,\rho^{-\alpha}
\qquad \forall \rho \geq \rho_*.
\]
Integrate from $\rho_*$ to $R > \rho_*$:
\[
\int_{\rho_*}^R \EE(\rho)^{-\beta}\,\EE'(\rho)\,d\rho
\geq c \int_{\rho_*}^R \rho^{-\alpha}\,d\rho.
\]
The left-hand side computes (using $\beta > 1$, so $1 - \beta < 0$):
\[
\int_{\rho_*}^R \EE^{-\beta}\,\EE'\,d\rho
= \Bigl[\frac{\EE^{1-\beta}}{1-\beta}\Bigr]_{\rho_*}^{R}
= \frac{1}{\beta - 1}\bigl(\EE(\rho_*)^{1 - \beta} - \EE(R)^{1-\beta}\bigr)
\leq \frac{\EE(\rho_*)^{1-\beta}}{\beta - 1},
\]
bounded uniformly in $R$ because $\EE(R)^{1-\beta} > 0$.

The right-hand side, using $\alpha < 1$:
\[
c\int_{\rho_*}^R \rho^{-\alpha}\,d\rho
= \frac{c}{1 - \alpha}\bigl(R^{1-\alpha} - \rho_*^{1-\alpha}\bigr)
\xrightarrow{R \to \infty} +\infty.
\]

Comparing: a bounded quantity exceeds an unbounded one, which is
impossible. The contradiction rules out
$\EE(\rho) \to \infty$. Therefore $\EE$ is bounded and
$\int_{\RR^3} \abs{\nabla v}^2 dx < \infty$.
\end{proof}

\begin{remark}[Quantitative form of the bound]
\label{rem:quantitative}
The proof of Theorem~\ref{thm:finite-D} yields an explicit
quantitative bound on $\EE$: from $\frac{1}{\beta -
1}\EE(\rho_*)^{1 - \beta} \geq \frac{c}{1-\alpha}(R^{1-\alpha} -
\rho_*^{1-\alpha})$ taken in the contrapositive direction, any
large value of $\EE(\rho_*)$ is constrained by the divergence
rate of $R^{1-\alpha}$. This gives the scale estimate
$\EE(\rho_0) \leq O(M^6/\rho_0)$ exploited in
Corollary~\ref{cor:quant}; see \S\ref{sec:extensions}.
\end{remark}

\section{From finite Dirichlet energy to $v \equiv 0$}
\label{sec:vzero}

We close the proof of Theorem~\ref{thm:main}. The finite
Dirichlet energy from Theorem~\ref{thm:finite-D}, combined with
$v \in \Lweak{3}$, upgrades $v$ to global $L^6$ regularity and
the pressure $p$ to global $L^3$; a global cutoff test then
forces $\nabla v \equiv 0$, hence $v \equiv 0$.

\begin{lemma}[Upgraded regularity]\label{lem:upgrade}
From $v \in \Lweak{3}(\RR^3)$ and $\int_{\RR^3} \abs{\nabla v}^2
dx < \infty$, it follows that $v \in L^6(\RR^3)$ with
$\norm{v}_{L^6} \leq C \norm{\nabla v}_{L^2}$, and the pressure
$p$ associated with the stationary NS system satisfies $p \in
L^3(\RR^3)$ with $\norm{p}_{L^3} \leq C\norm{\nabla v}_{L^2}^2$.
\end{lemma}

\begin{proof}
The homogeneous Sobolev inequality on $\RR^3$~\cite[II.5]{Galdi}
gives
\[
\norm{v}_{L^6(\RR^3)} \leq C_{\mathrm{Sob}}\,\norm{\nabla v}_{L^2(\RR^3)}.
\]
(Strictly: $v$ is assumed smooth, so Sobolev applies pointwise to
$v - v_R$ on any $B_R$ with $v_R$ the mean, and passing $R \to
\infty$ uses $v \in \Lweak{3}$ to eliminate the constant.) Hence
$v \otimes v \in L^3(\RR^3)$ with $\norm{v \otimes v}_{L^3} =
\norm{v}_{L^6}^2 \leq C\norm{\nabla v}_{L^2}^2$.

The stationary NS equation $-\nu \Delta v + \dv(v \otimes v)
+ \nabla p = 0$ gives, taking divergence and using $\dv v = 0$,
\[
-\Delta p = \dv \dv (v \otimes v) = \partial_i \partial_j (v_i v_j).
\]
In terms of Riesz transforms $R_i = (-\Delta)^{-1/2}\partial_i$,
$p = R_i R_j (v_i v_j)$ modulo the additive constant, which is
fixed to zero by requiring $p \in L^3$ (no nonzero constant lies
in $L^3(\RR^3)$). Calder\'on--Zygmund boundedness of $R_iR_j$ on
$L^3(\RR^3)$ gives $\norm{p}_{L^3} \leq C\norm{v\otimes v}_{L^3}
\leq C\norm{\nabla v}_{L^2}^2$.
\end{proof}

\begin{proof}[Proof of Theorem~\ref{thm:main}]
Let $\varphi_R \in C^\infty_c(\RR^3)$ be a radial cutoff with
$\varphi_R \equiv 1$ on $B_R$, $\supp \varphi_R \subset B_{2R}$,
$\abs{\nabla \varphi_R} \leq C/R$, and $\supp \nabla\varphi_R
\subset A(R, 2R)$. Test the stationary NS equation against
$v\,\varphi_R^2$ and integrate over $\RR^3$:
\[
\int \Bigl[-\nu\,\Delta v + (v\cdot\nabla)v + \nabla p\Bigr] \cdot v\,\varphi_R^2 \,dx = 0.
\]
Each term integrates by parts (valid by $v$ smooth, $\varphi_R$
compactly supported):
\begin{align*}
-\nu \int (\Delta v) \cdot v\,\varphi_R^2
&= \nu \int \abs{\nabla v}^2 \varphi_R^2 + \nu \int \varphi_R\,(\nabla\abs{v}^2) \cdot \nabla\varphi_R, \\
\int (v\cdot\nabla v) \cdot v\,\varphi_R^2
&= -\int \abs{v}^2\,(v \cdot \nabla\varphi_R)\,\varphi_R, \\
\int \nabla p \cdot v\,\varphi_R^2
&= -2\int p\,(v\cdot\nabla\varphi_R)\,\varphi_R,
\end{align*}
the last two using $\dv v = 0$. Rearranging gives the identity
\begin{equation}\label{eq:closing}
\nu \int \abs{\nabla v}^2 \varphi_R^2\, dx
= -\nu\int \varphi_R\,(\nabla\abs{v}^2)\cdot\nabla\varphi_R\,dx
+ \int \abs{v}^2\,(v\cdot\nabla\varphi_R)\,\varphi_R\, dx
+ 2\int p\,(v\cdot\nabla\varphi_R)\,\varphi_R\,dx.
\end{equation}

Denote the three right-hand terms $J_1, J_2, J_3$. Each is
supported in $A(R, 2R)$.

\emph{Bound on $J_1$ (viscous cross-term).} By Cauchy--Schwarz
and $\abs{\nabla \varphi_R} \leq C/R$,
\[
\abs{J_1} = 2\nu \abs{\int \varphi_R\,(v \cdot \nabla v)\cdot\nabla\varphi_R\,dx}
\leq \frac{2\nu C}{R}\,\norm{v}_{L^2(A(R,2R))}\,\norm{\nabla v}_{L^2(A(R, 2R))}.
\]
By Corollary~\ref{cor:lorentz-v} at $r = 2$,
$\norm{v}_{L^2(A(R,2R))} \leq C M R^{1/2}$. Since
$\int_{\RR^3}\abs{\nabla v}^2 < \infty$,
$\norm{\nabla v}_{L^2(A(R,2R))}^2 \leq \int_{\abs{x} \geq R}
\abs{\nabla v}^2\,dx \to 0$ as $R \to \infty$. Hence
\[
\abs{J_1} \leq \frac{C\nu M}{R^{1/2}}\,\norm{\nabla v}_{L^2(A(R,2R))}
\xrightarrow{R\to\infty} 0.
\]

\emph{Bound on $J_2$ (convective).} Using
$\abs{\nabla\varphi_R} \leq C/R$,
\[
\abs{J_2} \leq \frac{C}{R}\,\int_{A(R, 2R)} \abs{v}^3\,dx
= \frac{C}{R}\,\norm{v}_{L^3(A(R, 2R))}^3.
\]
By Proposition~\ref{prop:GN-critical} applied on
$A(R, 2R)$ (i.e., with $\theta = 1/2$ and $\rho = 2R$),
\[
\norm{v}_{L^3(A)}^3 \leq C'_{\mathrm{GN}}\bigl(M^a\,(2R)^b\,\DD_R^b + M^3\bigr),
\qquad \DD_R := \norm{\nabla v}_{L^2(A(R, 2R))}^2,
\]
with $a = 3/2$, $b = 3/4$ at $p = 2$. So
\[
\abs{J_2} \leq \frac{C}{R}\bigl(CM^{3/2}(2R)^{3/4}\DD_R^{3/4} + CM^3\bigr)
= C M^{3/2} R^{-1/4}\,\DD_R^{3/4} + C M^3 R^{-1} \xrightarrow{R\to\infty} 0,
\]
using $\DD_R \to 0$ and $R^{-1} \to 0$.

\emph{Bound on $J_3$ (pressure).} Using $\abs{\nabla\varphi_R}
\leq C/R$ and H\"older with $(3, 3/2)$,
\[
\abs{J_3} \leq \frac{2C}{R}\int_{A(R, 2R)} \abs{p}\,\abs{v}\,dx
\leq \frac{2C}{R}\,\norm{p}_{L^3(A(R, 2R))}\,\norm{v}_{L^{3/2}(A(R, 2R))}.
\]
By Corollary~\ref{cor:lorentz-v} at $r = 3/2$,
$\norm{v}_{L^{3/2}(A(R, 2R))} \leq CMR$. And by
Lemma~\ref{lem:upgrade}, $p \in L^3(\RR^3)$, so
$\norm{p}_{L^3(A(R, 2R))} \leq \norm{p}_{L^3(\abs{x}\geq R)}
\to 0$ as $R \to \infty$ by absolute continuity of the $L^3$
integral. Hence
\[
\abs{J_3} \leq \frac{2C \cdot CMR}{R}\,\norm{p}_{L^3(\abs{x}\geq R)}
= 2C^2M\,\norm{p}_{L^3(\abs{x}\geq R)} \xrightarrow{R\to\infty} 0.
\]

\emph{Passage to the limit.} From~\eqref{eq:closing},
\[
\nu \int \abs{\nabla v}^2 \varphi_R^2\,dx = J_1 + J_2 + J_3,
\]
and the left-hand side converges to $\nu \int_{\RR^3}
\abs{\nabla v}^2 dx$ as $R \to \infty$ by monotone convergence
(since $\varphi_R^2 \to 1$ pointwise, $0 \leq \varphi_R^2 \leq
1$, $\abs{\nabla v}^2 \in L^1$). The right-hand side converges to
$0$. Thus
\[
\nu \int_{\RR^3} \abs{\nabla v}^2\, dx = 0,
\qquad \text{hence } \nabla v \equiv 0,
\qquad \text{hence } v \equiv \text{const.}
\]
A nonzero constant does not lie in $\Lweak{3}(\RR^3)$: the
distribution function $\abs{\{x \in \RR^3 : \abs{v(x)} > t\}}$
is $\infty$ for $0 < t < \abs{v}$ and $0$ for $t > \abs{v}$, so
$\sup_t t\,\abs{\{\abs{v}>t\}}^{1/3} = \infty$. By hypothesis
$\norm{v}_{\Lweak{3}} < \infty$, so the constant is zero: $v
\equiv 0$. \qedhere
\end{proof}

\section{Extensions}\label{sec:extensions}

\subsection{Four-dimensional Liouville.}
\label{ssec:4d}

The proof of Theorem~\ref{thm:main} transposes to $\RR^n$ with
$n \geq 3$ by tracking the dimensional discriminant. The
Gagliardo--Nirenberg scale-invariance computation of
\S\ref{sec:gn} gives, in dimension $n$ at critical exponent
$L^n$,
\[
b = b(n, p) = \frac{n(n - p)}{2n - p(n - 2)},
\qquad a(n, p) = n - 2b.
\]
Substituting into the ODE exponents,
\[
\gamma(n, p) = -\frac{(n - 4)(n - 2p)}{2n - p(n - 2)},
\qquad
\alpha(n, p) = 1 + \frac{\gamma(n,p)}{b(n,p)},
\qquad
\beta(n, p) = \frac{1}{b(n, p)}.
\]
At $n = 3$ these reduce to $\gamma(3,p) = (3-2p)/(6-p)$ and
$b(3,p) = 3(3-p)/(6-p)$, recovering the exponents of
Theorem~\ref{thm:main-estimate} and Proposition~\ref{prop:ODE}.
The factor $(n - 4)$ in the numerator of $\gamma(n,p)$ governs
the transition across dimension four.

\begin{proof}[Proof of Corollary~\ref{cor:4d}]
At $n = 4$, $\gamma(4, p) = 0$ for every $p$; the GN main term
is scale-invariant and no $\rho$-weight survives. Taking the
natural choice $p = 5/2$ (symmetric about the critical
$L^{4,\infty}$): $b = 4 \cdot (3/2) / (8 - 5/2 \cdot 2) =
4 \cdot (3/2) / 3 = 2$; $\alpha = 1$; $\beta = 1/2$. This gives
a borderline ODE $\EE'(\rho) \geq c\,\rho^{-1}\,\EE(\rho)^{\beta}$
with $\alpha = 1$, which is integrable on the left (bounded
antiderivative since $\beta > 1$ --- for $p$ slightly larger than
$5/2$, say $p \in (5/2, 4)$, one obtains $\beta > 1$ by direct
substitution) and logarithmically divergent on the right:
\[
c\int_{\rho_*}^R \frac{d\rho}{\rho}
= c\log\frac{R}{\rho_*}
\xrightarrow{R \to \infty} +\infty.
\]
The contradiction closes the same way as in
Theorem~\ref{thm:finite-D}, forcing $\int_{\RR^4}\abs{\nabla v}^2 dx
< \infty$. The closing argument of \S\ref{sec:vzero} extends
verbatim with Sobolev $\norm{v}_{L^4(\RR^4)} \leq C\norm{\nabla v
}_{L^2(\RR^4)}$ (critical Sobolev in dimension four) replacing
the three-dimensional $L^6$ bound. Hence $v \equiv 0$ on $\RR^4$.
\end{proof}

\begin{remark}[Failure for $n \geq 5$]
\label{rem:high-dim}
For $n \geq 5$, $\gamma(n, p) > 0$ for $p$ in the admissible
range, making $\alpha(n, p) > 1$. The integral $\int \rho^{-\alpha} d\rho$
converges, and the RHS of the ODE integration stays bounded --- no
contradiction. A sharper structural obstruction applies: $\beta
> 1$ requires $p > n/2$ (from the Sobolev embedding direction),
while $\alpha \leq 1$ requires $p \leq n/2$ (from the integration
direction); the two conditions are incompatible for $n \geq 5$.
No stationary $L^{n,\infty}(\RR^n)$ Liouville theorem
is currently available for $n \geq 5$; the structural obstruction
(Sobolev / integration exponent conflict) is discussed in the
companion corpus \cite{CompanionCorpus}.
\end{remark}

\subsection{The $W^{1,2}_\loc$ weak-solution extension.}
\label{ssec:weak}

\begin{proof}[Proof of Corollary~\ref{cor:weak}]
Let $v \in \Lweak{3}(\RR^3) \cap W^{1,2}_\loc(\RR^3)$ be a weak
solution of stationary Navier--Stokes in the sense that, for every
divergence-free $\phi \in C^\infty_c(\RR^3)$,
\[
\nu \int \nabla v : \nabla \phi\, dx
+ \int v \otimes v : \nabla \phi\, dx = 0.
\]
(The pressure is eliminated by the divergence-free test class.)

\emph{Step 1: local regularity.} By $W^{1,2}_\loc$ and the
Sobolev embedding $W^{1,2}(B_R) \hookrightarrow L^6(B_R)$,
$v \in L^6_\loc$. Interpolating
$v \in \Lweak{3} \cap L^6_\loc$ with H\"older gives
$v \in L^r_\loc$ for every $r \in [3, 6]$. In particular
$v \otimes v \in L^3_\loc$, so $p \in L^3_\loc$ by the local
Calder\'on--Zygmund argument of Lemma~\ref{lem:upgrade} restricted
to any $B_R$. Iteration with the equation $-\nu\Delta v + \dv(v
\otimes v) + \nabla p = 0$ (in the sense of distributions) gives
$v \in C^\infty_\loc$ (Galdi~\cite[Ch.~V]{Galdi}).

\emph{Step 2: the Caccioppoli and GN machinery still apply.} The
test-function construction of \S\ref{sec:pressure} uses
$w = v\eta - \BB[v\cdot\nabla\eta]$; this requires only $v \in
W^{1,q}_\loc$ for some $q > 1$, which follows from the interior
regularity of Step~1. Proposition~\ref{prop:elim} holds by
mollification of $v$ followed by passage to the limit, since
$v$ is smooth in the interior. The Lorentz and GN machinery of
\S\ref{sec:lorentz}--\S\ref{sec:gn} use $v \in \Lweak{3}$ and
local Sobolev regularity, both of which persist under the weak
hypothesis.

\emph{Step 3: test-function density.} The divergence-free test
space $\{\phi \in C^\infty_c(\RR^3) : \dv \phi = 0\}$ is dense
in the space $W^{1,q}_0$-divergence-free functions with compact
support, for every $q > 1$~\cite[Thm~III.4.1]{Galdi} (originally
due to Heywood~\cite{Heywood}). The Bogovski\u{\i}-modified test
function $w = v\eta - \BB[v\cdot\nabla\eta]$ lies in this space
by construction, so it is an admissible test function in the
weak formulation.

\emph{Step 4: closing argument.} The proof of
Theorem~\ref{thm:main} runs verbatim, concluding $v \equiv 0$.
The $W^{1,2}_\loc$ hypothesis is sharp: the Landau
solutions~\cite{Landau} are smooth outside the origin, belong to
$\Lweak{3}(\RR^3)$ (with quasi-norm equal to a definite constant),
and lie in $W^{1,2}_\loc(\RR^3 \setminus \{0\})$ but not at the
origin; they show that \emph{any weakening below $W^{1,2}_\loc$
admits nontrivial stationary solutions} in $\Lweak{3}$.
\end{proof}

\subsection{Quantitative bound.}
\label{ssec:quantitative}

\begin{proof}[Proof of Corollary~\ref{cor:quant}]
From the proof of Proposition~\ref{prop:ODE}, the ODE constant
at $p = 2$ is
\[
c = \frac{\theta}{(1 - \theta)(2K_2)^{1/b}}
= \frac{\theta}{(1 - \theta)(2K_2)^{4/3}},
\]
with $K_2 = O(M^{3/2})$ (tracking $M$-dependence through
Theorem~\ref{thm:main-estimate}; the dominant $M$-power arises
from $\norm{v}_{L^p}^a$ with $a = 3/2$ at $p = 2$). So $c \geq
c_0 M^{-2}$ for a dimensionless $c_0 = c_0(\theta, \nu)$.

From the integration of
Theorem~\ref{thm:finite-D},
\[
\frac{\EE(\rho_*)^{1 - \beta}}{\beta - 1}
\geq \frac{c}{1 - \alpha}\bigl(R^{1-\alpha} - \rho_*^{1-\alpha}\bigr)
\qquad\forall R > \rho_*.
\]
At $p = 2$: $\beta - 1 = 1/3$, $1 - \alpha = 1/3$. Taking
$R = 2\rho_*$,
\[
\EE(\rho_*)^{-1/3} \geq \frac{c \cdot (2^{1/3} - 1)}{1/3 \cdot 1/3}\,
\rho_*^{1/3}
= C_0\,c\,\rho_*^{1/3}.
\]
Inverting,
\[
\EE(\rho_*) \leq \frac{C_1}{c^3\,\rho_*}
= \frac{C_1 M^6}{c_0^3\,\rho_*},
\]
which is the claimed single-exponential-in-$M$ bound.

\emph{Comparison with Tao.}
Tao's quantitative regularity bound~\cite{Tao2019}, which targets
the full Clay problem via an assumed uniform $L^3$-bound on the
evolution, gives a triple-exponential dependence on the bound.
Corollary~\ref{cor:quant} is, in principle, not directly
comparable (it concerns the stationary case at a single scale),
but the constructivized KNSS reduction using
Theorem~\ref{thm:main} as the Liouville step in place of a
conjectured one would inherit the present polynomial-in-$M$ bound.
\end{proof}

\section{Application to the Clay chain and the $L^{3,\infty}$ endpoint gap}
\label{sec:clay-chain}

This section describes how Theorem~\ref{thm:main} composes within
the Koch--Nadirashvili--Seregin--\v{S}ver\'ak blow-up framework,
identifies the single structural step that remains open, and
surveys the current partial closures.

\subsection{The composition.}
Suppose $u \in C^\infty(\RR^3 \times [0, T^*))$ is the maximal
smooth solution of~\eqref{eq:NS} with Schwartz $u_0$, and
$T^* < \infty$. The chain runs:
\begin{enumerate}[label=(C\arabic*),itemsep=2pt]
\item \textbf{Type I excluded.} NRS~\cite{NRS96},
Tsai~\cite{Tsai98}, ESS~\cite{ESS03RMS}: no self-similar
solutions; any hypothetical singularity is Type II.
\item \textbf{KNSS rescaling.}~\cite[\S6]{KNSS}: sequences
$\lambda_n \to 0$, $x_n \to x_0$, $t_n \to T^*$ with
$v_n(y, s) := \lambda_n u(x_n + \lambda_n y, t_n + \lambda_n^2 s)
\to v$ in $C^m_\loc$, $\abs{v} \leq C_0$,
$\abs{v(0, 0)} = 1$.
\item \textbf{Stationary backward limit.} Arzel\`a--Ascoli
extracts $v(\cdot, s_n) \to v_\infty$ in $C^m_\loc$ along some
$s_n \to -\infty$; equicontinuity of $\partial_s v$ forces
$v_\infty$ stationary.
\item \textbf{Liouville applies.} Under the inheritance step
discussed in \S\ref{ssec:inheritance}, $v_\infty \in
\Lweak{3}$; by Theorem~\ref{thm:main}, $v_\infty \equiv 0$.
\item \textbf{Backward uniqueness contradicts.}
ESS~\cite[\S5]{ESS03RMS} + spatial continuation
\cite[\S4]{ESS03RMS} yield $\omega(\cdot, s) \neq 0$ for all
$s \leq 0$, contradicting the full convergence from (C4).
\end{enumerate}

Every step of this chain except (C4) is rigorous in the
published literature. The exception is the implicit hypothesis
that $v \in \Lweak{3}$ uniformly in $s \leq 0$; the conclusion
of Theorem~\ref{thm:main} is conditional on that uniform bound.

\subsection{The inheritance gap.}\label{ssec:inheritance}

Scale invariance of the $\Lweak{3}$ quasi-norm shows that the
ancient solution $v$ inherits $\norm{v(\cdot, s)}_{\Lweak{3}}$
from $\norm{u(\cdot, t_n + \lambda_n^2 s)}_{\Lweak{3}}$. What
remains is to bound the latter uniformly in $t \in [0, T^*)$.
This is the question of the $L^{3,\infty}$ endpoint regularity
criterion, which Phuc~\cite[p.~4]{Phuc2015} states as open:
\begin{quote}
``To the best of our knowledge, a criterion of local regularity
for the Navier--Stokes equations in $L^\infty(-1, 0;
L^{3,\infty}(B_1))$ is still unknown.''
\end{quote}

The obstruction is structural. Phuc's method closes the
criterion for $L^\infty_t L^{3,q}_x$ with $q \in [3, \infty)$
using exterior-tail compactness at the blow-up extraction: the
rescaled ancient solution inherits $\norm{\cdot}_{L^{3,q}
(\RR^3 \setminus B_R)} \to 0$ as $R \to \infty$ by absolute
continuity of the $L^{3,q}$ norm. At $q = \infty$, the
$\Lweak{3}$ quasi-norm is not absolutely continuous --- the
function $\abs{x}^{-1}$, the canonical critical profile, has
scale-invariant exterior norm --- and the compactness step
collapses.

\subsection{Partial closures.}

\paragraph{Axisymmetric case.}
O\.{z}a\'nski--Palasek~\cite{OzanskiPalasek} closed the endpoint
for axisymmetric flows via Nazarov--Ural'tseva parabolic Harnack
\cite{NazarovUraltseva} applied to the swirl $rv_\theta$, which
satisfies a scalar parabolic equation with divergence-free drift.
The reduction to a scalar quantity is essentially axisymmetric;
no analog is known for general flows.

\paragraph{Smallness.}
Luo--Tsai~\cite{LuoTsai} and Kim--Kozono~\cite{KimKozono} close
the endpoint under a Serrin-type smallness assumption
$\norm{u}_{L^\infty_t L^{3,\infty}_x}$ below a threshold. The
large-data question remains open.

\paragraph{Partial regularity.}
Seregin~\cite{Seregin2019} established that an
$L^\infty_t L^{3,\infty}_x$ bound controls the singular set to be
finite (bounded by $M^3/\epsilon^4$ for an effective $\epsilon$);
this is partial, not full, regularity.

\paragraph{Logarithmic improvement.}
Phuc~\cite[Thm~1.5]{Phuc2015} closes the endpoint under the
pointwise refinement
$\abs{u(x, t)} \leq C\abs{x}^{-1}\abs{\log(\abs{x}/2)}^{-\beta}$
for any $\beta > 0$. A strengthening of the rescaled
ancient-solution bound to this logarithmic form would complete
the chain via Phuc's machinery.

\subsection{A weighted-moment proposal.}\label{ssec:weighted-moment}

The obstruction at the $L^2$-only level is sharp: the
translation-invariance of $\norm{\cdot}_{L^2(\RR^3)}$ forbids any
constant depending only on $\norm{u_0}_{L^2}$ from bounding
$\abs{u(x, t)}$ pointwise. Brandolese--Meyer~\cite{BrandoleseMeyer}
exhibit explicit counterexamples. The constant must depend on a
weighted norm carrying spatial information.

A natural proposal is the weighted hypothesis
$(1 + \abs{x}) u_0 \in L^\infty(\RR^3)$, which is automatic for
Schwartz data. Define $\mathcal{M}(u_0) := \norm{u_0}_{L^2} +
\norm{(1 + \abs{x}) u_0}_{L^\infty}$. We conjecture that the
rescaled ancient solution inherits
$\norm{v(\cdot, s)}_{\Lweak{3}} \leq M(\mathcal{M}(u_0), \nu)$
uniformly, via a three-step strategy: (i) local weighted
$L^\infty$ propagation from Kato mild-solution theory + the
Brandolese--Meyer generic-decay framework; (ii) global
propagation up to $T^*$ via a Gronwall-type estimate on the
weighted quantity $w = (1 + \abs{x}) u$; (iii) scale-invariance
transfer to the rescaled sequence.

Step (iii) is elementary; step (i) is classical. Step (ii) is
the novel mathematical content, and closing it rigorously would
resolve the Clay regularity question for Schwartz data under this
mildly strengthened hypothesis. A preliminary estimate analysis
\cite[desk/04]{CompanionCorpus} shows that the naive bootstrap
argument gives a short-time bound (closed under a quadratic
Volterra inequality $\tilde A \leq L_0 + L_1 \tilde A^2$) but
does not obviously extend to all of $[0, T^*)$; the extension
requires either (a) a uniform-in-$t$ Caffarelli--Kohn--Nirenberg
$\epsilon$-regularity bound on the exterior of the singular set,
giving $\abs{\nabla^k u(x, t)} \leq C_k(1 + \abs{x})^{-1-k}$ for
$\abs{x} \geq R_0$, or (b) a weighted $L^2$ energy estimate that
propagates $\norm{(1 + \abs{x})^\alpha u(\cdot, t)}_{L^2}$ up to
$T^*$.

The full argument, together with the exterior
$\epsilon$-regularity analysis, is deferred to a forthcoming
work. The present paper establishes Theorem~\ref{thm:main}
unconditionally and records this closure as the open problem
reducing the Clay chain.

\section{The obstruction landscape}\label{sec:obstructions}

The research program connecting Theorem~\ref{thm:main} to the
bounded-Liouville conjecture for the full regularity problem
runs through nine structural obstructions catalogued in the
companion corpus \cite{CompanionCorpus}. We list them briefly
here; full statements, canonical citations, and alternative-route
discussions are in the corpus.

\begin{enumerate}[label=(O\arabic*), itemsep=2pt]
\item \textbf{Closure Obstruction}: no finite combination of the
fifty-six antecedent algebraic identities plus
Cauchy--Schwarz/Hodge/Helmholtz/coarea closes the Final
Sub-Question strain-$L^2$ improvement.
\item \textbf{Exchange-of-Limits Obstruction}: the hull topology
does not commute with the far-field Rayleigh-quotient limit
along translation-hull sequences.
\item \textbf{Riesz Obstruction}: the pressure Poisson equation
plus $L^2$ bounds does not determine the sign of the pressure
Hessian at the eigenfunction supremum.
\item \textbf{USC Obstruction}: upper semi-continuity of the
asymptotic Rayleigh quotient at constants is equivalent to, not
independent of, bounded Liouville.
\item \textbf{Carleman Obstruction}: unique continuation (UCP)
is categorically distinct from Liouville (non-existence); the
two cannot compose without exponent loss.
\item \textbf{Integral-Equation Obstruction}: under linear
hypotheses, the Birman--Schwinger integral equation admits
bounded eigenfunctions at negative eigenvalue; linear spectral
analysis is exhausted.
\item \textbf{Resonance-Diagonal Obstruction}: in the $\Phi$-adic
cross-shell coupling matrix on Fibonacci shells, the pressure
correction sign is undetermined per shell --- a per-scale version
of the Riesz obstruction.
\item \textbf{Helicity-Sector Obstruction}: the symmetric
geometric factor $\Lambda_k^{(S)}$ is independent of the scalar
quadratic NS invariants; helicity-based closures do not fix its
sign.
\item \textbf{Twist-Writhe Obstruction}: the
C\u{a}lug\u{a}reanu--White--Fuller decomposition and the sector
decomposition live in incommensurable representations; the
$\ell = 2$ angular harmonic on $\Phi$-shells is the irreducible
obstruction.
\end{enumerate}

The landscape is \emph{topological, not metric}: each obstruction
identifies a signed quantity with $O(C(M))$-controlled magnitude
but undetermined sign. Alternative-route research programs ---
hull-ergodic closure via the Grand Unification Conjecture,
icosahedral averaging of the $\ell = 2$ component on
$\Phi$-shells, space-time hull closure for Type II --- are
developed in the companion corpus.

\section{Discussion}\label{sec:discussion}

\subsection{The vortex-stretching reading.}

Theorem~\ref{thm:main} admits a physical reading. The weighted
energy $\EE(\rho) = \int \abs{\nabla v}^2 \varphi(\abs{x}/\rho)
dx$ measures \emph{fluid enstrophy} at scale $\rho$, and the
derivative $\EE'(\rho)$ is a local influx rate across the sphere
$\abs{x} = \rho$. The superlinear ODE
$\EE' \geq c\rho^{-\alpha}\EE^\beta$ with $\beta > 1$ states that
the influx rate grows superlinearly with the enstrophy already
present: vortex-stretching amplifies strain, which amplifies
vorticity, in a feedback loop. The argument is in spirit a
no-go: such a feedback, applied across the bulk of $\RR^3$,
cannot stabilize; the only fixed point is $v \equiv 0$.

The Lorentz hypothesis $v \in \Lweak{3}$ enters as the
sub-scaling-critical bound on the tail of $v$ that keeps the
feedback from saturating. At the endpoint $L^{3,\infty}$ the
bound is still strong enough, but only barely --- the $\rho^\gamma$
factor in the ODE degenerates at $p \to 3$, and vanishes at the
intermediate $p = 3/2$. Between these, the structural exponents
$\alpha < 1 < \beta$ are strict, and the argument closes.

\subsection{Geometric architecture.}

The proof relies on two distinct geometric choices that are
sharp:

\emph{Polynomial (fixed-ratio) annuli versus exponential weights.}
The fixed-ratio annulus $A(\theta\rho, \rho)$ with $\theta \in
(0, 1)$ fixed is geometrically scale-invariant, which is what
allows the Bogovski\u{\i} operator on $A(\theta\rho, \rho)$ to
carry uniform-in-$\rho$ constants (Corollary
\ref{cor:bogo-annulus}). A Gaussian or exponential weight
(common in the dynamical arguments of~\cite{KNSS, ESS03RMS})
gives strictly better decay on the tail but destroys
scale-invariance of the Bogovski\u{\i} constant, requiring
$\rho$-dependent bookkeeping that ultimately fails at the
critical exponent.

\emph{The hat function versus smooth cutoffs.} The sharp
hat-function $\varphi$ of Notation~\ref{not:geometry} gives the
exact derivative identity \eqref{eq:hat-deriv} with no
$\epsilon$-bookkeeping. A smoothed cutoff would require a
transition-layer argument that introduces additional surface
integrals; these wash out in the limit but obscure the structural
form of the inequality.

\subsection{Scope and limitations.}

Theorem~\ref{thm:main} treats the stationary case exclusively;
the evolution case is the KNSS composition discussed in
\S\ref{sec:clay-chain}, where the remaining bottleneck is the
uniform $L^{3,\infty}$ inheritance of the rescaled ancient
solution. The Liouville step itself is proved unconditionally.

The $n \geq 5$ case is inaccessible by the present method (see
Remark~\ref{rem:high-dim}). This is not an artefact: no
$L^{n,\infty}$ Liouville theorem is known for stationary NS on
$\RR^n$ with $n \geq 5$, and the structural obstruction (the
Sobolev / integration exponent conflict) suggests that different
machinery is needed.

The constants tracked in Corollary~\ref{cor:quant} are polynomial
in the initial $L^{3,\infty}$ bound $M$, which is a quantitative
improvement over conjectural triple-exponential bounds
(\cite{Tao2019}) that would arise from the full-evolution case.
However, the constants here are interior-to-the-argument; making
them optimal (in the $\theta, \nu$ dependence) is out of scope.

\subsection{Open problems.}

We close with four problems, in decreasing order of specificity.

\begin{enumerate}[label=(P\arabic*), itemsep=2pt]
\item \textbf{The $L^{3,\infty}$ endpoint.} Resolve Phuc's open
problem (\cite{Phuc2015}, reproduced in \S\ref{ssec:inheritance}):
is $\norm{u}_{L^\infty_t L^{3,\infty}_x(B_1 \times (-1, 0))}$ a
regularity criterion? Theorem~\ref{thm:main} is the stationary
half; the parabolic analog with backward uniqueness substitutes
is the outstanding mathematical question.
\item \textbf{The weighted-moment proposal.} Close the
weighted-moment closure of \S\ref{ssec:weighted-moment} for
Schwartz data: prove that $(1 + \abs{x})u_0 \in L^\infty(\RR^3)$
implies a uniform $L^{3,\infty}$ bound on the rescaled ancient
solution. The main open step is the exterior
$\epsilon$-regularity transport to obtain uniform $\abs{\nabla^k
u(x, t)} \leq C_k (1 + \abs{x})^{-1-k}$ for $\abs{x}$ large.
\item \textbf{Non-axisymmetric Nazarov--Ural'tseva.} Find a
non-axisymmetric scalar invariant of a 3D NS flow that satisfies
a parabolic equation with divergence-free drift amenable to
Nazarov--Ural'tseva Harnack~\cite{NazarovUraltseva}. The
axisymmetric $rv_\theta$ has no known general-case analog;
identifying one would generalize~\cite{OzanskiPalasek}.
\item \textbf{Lean formalization.} Formalize
Theorem~\ref{thm:main} in Lean 4/Mathlib. The companion corpus
\cite{CompanionCorpus} contains a skeleton with three open
\texttt{sorry}s in the Caccioppoli, ODE contradiction, and
Dirichlet-integral-zero steps. Closing those would upgrade the
present paper to machine-verified status.
\end{enumerate}

\subsection{Acknowledgments.}

I thank the Koch--Nadirashvili--Seregin--\v{S}ver\'ak program
for the blow-up-extraction framework that made the present work
coherent, and the Phuc line of work on $L^{3,q}$ regularity
criteria for clarifying the endpoint obstruction. Conversations
with the Turbulent Waters team and the HIRA kernel project on
the obstruction landscape influenced \S\ref{sec:obstructions}.
The error-ledger methodology that produced CORR-05/08/09
\cite{CompanionCorpus} --- catching attribution errors before
publication --- is a collaborative discipline whose value
exceeded any single iteration's correctness.

\bigskip

% ============================================================
\begin{thebibliography}{99}

\bibitem{Bogovskii}
M.~E.~Bogovski\u{\i},
\emph{Solution of the first boundary value problem for the equation of continuity of an incompressible medium},
Dokl.\ Akad.\ Nauk SSSR \textbf{248} (1979), 1037--1040.

\bibitem{BrandoleseMeyer}
L.~Brandolese and Y.~Meyer,
\emph{On the instantaneous spreading for the Navier--Stokes system in the whole space},
ESAIM: Control, Optimisation and Calculus of Variations
\textbf{8} (2002), 273--285.

\bibitem{Brandolese2004}
L.~Brandolese,
\emph{Space-time decay of Navier--Stokes flows invariant under rotations},
Math.\ Ann.\ \textbf{329} (2004), 685--706.

\bibitem{CKN}
L.~Caffarelli, R.~Kohn, and L.~Nirenberg,
\emph{Partial regularity of suitable weak solutions of the Navier--Stokes equations},
Comm.\ Pure Appl.\ Math.\ \textbf{35} (1982), 771--831.

\bibitem{CNY2024}
Y.~Cho, J.~Neustupa, and M.~Yang,
\emph{A regularity criterion for the stationary Navier--Stokes equations in the Lorentz space $L^{3,\infty}(\RR^3)$},
Nonlinearity \textbf{37} (2024), Paper No.\ 035007.

\bibitem{CY2026}
Y.~Cho and M.~Yang,
\emph{Refined Liouville-type theorems for the stationary Navier--Stokes equations},
arXiv:2603.23833 (2026).

\bibitem{CompanionCorpus}
A.~Nahiru, \emph{finish\_navier\_stokes companion corpus}, GitHub
repository, \url{https://github.com/hirunation/finish_navier_stokes},
archived at Zenodo, DOI: 10.XXXX/zenodo.XXXXXXX (2026).

\bibitem{ESS03ARMA}
L.~Escauriaza, G.~Seregin, and V.~\v{S}ver\'ak,
\emph{Backward uniqueness for parabolic equations},
Arch.\ Ration.\ Mech.\ Anal.\ \textbf{169} (2003), 147--157.

\bibitem{ESS03RMS}
L.~Escauriaza, G.~Seregin, and V.~\v{S}ver\'ak,
\emph{$L^{3,\infty}$-solutions of Navier--Stokes equations and backward uniqueness},
Russ.\ Math.\ Surv.\ \textbf{58} (2003), 211--250.

\bibitem{Galdi}
G.~P.~Galdi,
\emph{An Introduction to the Mathematical Theory of the Navier--Stokes Equations},
Springer Monographs in Mathematics, 2nd ed., Springer, 2011.

\bibitem{Heywood}
J.~G.~Heywood,
\emph{The Navier--Stokes equations: on the existence, regularity and decay of solutions},
Acta Math.\ \textbf{136} (1976), 61--102.

\bibitem{Kato}
T.~Kato,
\emph{Strong $L^p$-solutions of the Navier--Stokes equation in $\RR^m$, with applications to weak solutions},
Math.\ Z.\ \textbf{187} (1984), 471--480.

\bibitem{KimKozono}
H.~Kim and H.~Kozono,
\emph{Interior regularity criteria in weak spaces for the Navier--Stokes equations},
Manuscripta Math.\ \textbf{115} (2004), 85--100.

\bibitem{KNSS}
G.~Koch, N.~Nadirashvili, G.~Seregin, and V.~\v{S}ver\'ak,
\emph{Liouville theorems for the Navier--Stokes equations and applications},
Acta Math.\ \textbf{203} (2009), 83--105.

\bibitem{Landau}
L.~D.~Landau,
\emph{On a new exact solution of the Navier--Stokes equations},
Dokl.\ Akad.\ Nauk SSSR \textbf{43} (1944), 299--301.

\bibitem{Leray}
J.~Leray,
\emph{Sur le mouvement d'un liquide visqueux emplissant l'espace},
Acta Math.\ \textbf{63} (1934), 193--248.

\bibitem{LuoTsai}
W.~Luo and T.-P.~Tsai,
\emph{Regularity criteria in weak $L^3$ for 3D incompressible Navier--Stokes equations},
Funkcialaj Ekvacioj \textbf{58} (2015), 387--404.

\bibitem{NRS96}
J.~Ne\v{c}as, M.~R\r{u}\v{z}i\v{c}ka, and V.~\v{S}ver\'ak,
\emph{On Leray's self-similar solutions of the Navier--Stokes equations},
Acta Math.\ \textbf{176} (1996), 283--294.

\bibitem{NazarovUraltseva}
A.~I.~Nazarov and N.~N.~Ural'tseva,
\emph{The Harnack inequality and related properties of solutions to elliptic and parabolic equations with divergence-free lower-order coefficients},
Algebra i Analiz \textbf{23} (2011), 136--168.

\bibitem{OzanskiPalasek}
W.~S.~O\.{z}a\'nski and S.~Palasek,
\emph{Quantitative regularity for axisymmetric Navier--Stokes solutions with a bound in the weak-$L^3$ critical norm},
arXiv:2210.10030, Annals of PDE (2023).

\bibitem{Palasek}
S.~Palasek,
\emph{Improved quantitative regularity for the Navier--Stokes equations in a scale of critical spaces},
Arch.\ Ration.\ Mech.\ Anal.\ \textbf{242} (2021).

\bibitem{Phuc2015}
N.~C.~Phuc,
\emph{The Navier--Stokes equations in nonendpoint borderline Lorentz spaces},
J.~Math.\ Fluid Mech.\ \textbf{17} (2015), 741--760.

\bibitem{Schonbek}
M.~E.~Schonbek,
\emph{$L^2$ decay for weak solutions of the Navier--Stokes equations},
Arch.\ Ration.\ Mech.\ Anal.\ \textbf{88} (1985), 209--222.

\bibitem{Seregin2019}
G.~Seregin,
\emph{A note on weak solutions to the Navier--Stokes equations that are locally in $L^\infty(L^{3,\infty})$},
arXiv:1906.06707 (2019).

\bibitem{Tao2019}
T.~Tao,
\emph{Quantitative bounds for critically bounded solutions to the Navier--Stokes equations},
Contemp.\ Math.\ \textbf{740} (2020), AMS.

\bibitem{Tsai98}
T.-P.~Tsai,
\emph{On Leray's self-similar solutions of the Navier--Stokes equations satisfying local energy estimates},
Arch.\ Ration.\ Mech.\ Anal.\ \textbf{143} (1998), 29--51.

\bibitem{Wiegner}
M.~Wiegner,
\emph{Decay results for weak solutions of the Navier--Stokes equations on $\RR^n$},
J.~London Math.\ Soc.\ \textbf{35} (1987), 303--313.

\end{thebibliography}

\end{document}
